\documentclass[11pt,a4paper]{article}

\usepackage[T1]{fontenc}
\usepackage[utf8]{inputenc}
\usepackage{lmodern}
\usepackage[margin=2.3cm]{geometry}
\usepackage{xcolor}
\usepackage{booktabs}
\usepackage{tabularx}
\usepackage{enumitem}
\usepackage{listings}
\usepackage[most]{tcolorbox}
\usepackage{amsmath,amssymb}
\usepackage{braket}
\usepackage{textcomp}
\usepackage{etoolbox}
% Greek letters and symbols typed directly in \texttt and \cmd arguments.
\DeclareUnicodeCharacter{0394}{\ensuremath{\Delta}}
\DeclareUnicodeCharacter{03A9}{\ensuremath{\Omega}}
\DeclareUnicodeCharacter{03C1}{\ensuremath{\rho}}
\DeclareUnicodeCharacter{03C0}{\ensuremath{\pi}}
\DeclareUnicodeCharacter{2248}{\ensuremath{\approx}}
\usepackage{xurl}
\usepackage[colorlinks=true,linkcolor=blue!50!black,urlcolor=blue!50!black,
  citecolor=blue!50!black]{hyperref}

% ---------------------------------------------------------------------------
% Edit these before sending the document out.
\newcommand{\coursedate}{\textbf{21.10.2026}}
\newcommand{\courserepourl}{\url{https://github.com/marcelojbp/physics_RSE_and_AI}}
% ---------------------------------------------------------------------------

\definecolor{codebg}{gray}{0.95}
\lstset{
  basicstyle=\ttfamily\small,
  backgroundcolor=\color{codebg},
  frame=single, rulecolor=\color{gray!40},
  breaklines=true, columns=fullflexible, keepspaces=true,
  showstringspaces=false, upquote=true, escapeinside={(*}{*)},
  xleftmargin=3pt, xrightmargin=3pt, aboveskip=6pt, belowskip=6pt,
  extendedchars=true, inputencoding=utf8,
  literate=
    {Δ}{{$\Delta$}}1 {Ω}{{$\Omega$}}1 {π}{{$\pi$}}1 {ρ}{{$\rho$}}1
    {σ}{{$\sigma$}}1 {θ}{{$\theta$}}1 {φ}{{$\varphi$}}1 {δ}{{$\delta$}}1
    {Γ}{{$\Gamma$}}1 {μ}{{$\mu$}}1 {ħ}{{$\hbar$}}1 {≈}{{$\approx$}}1
    {≤}{{$\leq$}}1 {≥}{{$\geq$}}1 {±}{{$\pm$}}1 {→}{{$\to$}}1
    {⟨}{{$\langle$}}1 {⟩}{{$\rangle$}}1 {₁}{{$_1$}}1 {₂}{{$_2$}}1
    {τ}{{$\tau$}}1 {ψ}{{$\psi$}}1 {ω}{{$\omega$}}1 {²}{{$^2$}}1 {⁴}{{$^4$}}1
    {½}{{$\frac12$}}1 {†}{{$^\dagger$}}1 {₊}{{$_+$}}1 {₋}{{$_-$}}1 {ₑ}{{$_e$}}1
    {−}{{$-$}}1 {√}{{$\surd$}}1 {∞}{{$\infty$}}1 {–}{{--}}1 {—}{{---}}1
}
% Files quoted from the Julia project, so that the notes stay in sync with it.
\newcommand{\ds}{../DrivenSpin}

\newtcolorbox{note}[1][Note]{colback=blue!4, colframe=blue!45!black,
  fonttitle=\bfseries, title=#1, boxrule=0.6pt, arc=2pt}
\newtcolorbox{warn}[1][Watch out]{colback=orange!6, colframe=orange!70!black,
  fonttitle=\bfseries, title=#1, boxrule=0.6pt, arc=2pt}
% Questions participants may ask, with short answers (mainly for the instructor).
\newtcolorbox{qa}[1][Questions you may get]{colback=gray!4, colframe=gray!55!black,
  fonttitle=\bfseries, title=#1, boxrule=0.6pt, arc=2pt}
% Placeholders: code from DrivenSpin/ that does not exist yet, text to be extended.
\newtcolorbox{placeholder}[1][Placeholder]{colback=red!4, colframe=red!60!black,
  fonttitle=\bfseries, title=#1, boxrule=0.6pt, arc=2pt}

\newcommand{\cmd}[1]{\texttt{#1}}
\newcommand{\onslides}[1]{\par\noindent\textit{On the slides:} #1\par}
\newcommand{\Q}[1]{\par\smallskip\noindent\textbf{Q: #1}\par\noindent}

% Physics (PLAN.md section 5)
\newcommand{\dd}{\mathrm{d}}
\newcommand{\ee}{\mathrm{e}}
\newcommand{\tr}{\mathrm{tr}}
\newcommand{\Pe}{P_e}
\newcommand{\OmR}{\Omega_{\mathrm{R}}}
\newcommand{\Dis}{\mathcal{D}}

% References: \rf{key} or \rf{key1,key2} link to entries of the reference list.
\newcommand{\rf}[1]{[\def\rfsep{}\forcsvlist{\rfone}{#1}]}
\newcommand{\rfone}[1]{\rfsep\ref{ref:#1}\def\rfsep{, }}

\setlist{itemsep=2pt, topsep=4pt}
\setlength{\parindent}{0pt}
\setlength{\parskip}{5pt}
\setlength{\emergencystretch}{2em}

\title{\textbf{AI for Physics Projects}\\[4pt]
\large Course notes: background, reference cards and annotated references}
\author{}
\date{Course of \coursedate}

\begin{document}
\maketitle
\vspace{-1.5em}

\begin{note}[How to use these notes]
These notes follow the slides block by block and add what does not fit on a slide:
the reasoning behind each claim, the details a sharp question will reach, and where
to read more. They serve three readers.
\begin{itemize}
  \item \textbf{The instructor}: each section explains the claims on its slides in
        enough depth to answer questions, and ends with a box of questions you may get.
  \item \textbf{Participants who want to go further}: Part~II has two reference cards
        (git and a verification checklist) that also work on their own; each section
        points to the best next reading.
  \item \textbf{Anyone checking a source}: Part~III lists every reference with a link,
        what it is good for, and how it was verified. Numbers in square brackets, such
        as \rf{tao}, point there.
\end{itemize}
\end{note}

\textbf{Conventions.} Throughout, $\hbar = 1$; times are in \textmu s and angular
frequencies in rad/\textmu s, so a drive at $1$\,MHz has angular frequency
$2\pi \times 1\,\mathrm{rad}/\mu\mathrm{s}$. The basis order is $(\ket{e}, \ket{g})$,
with $\ket{e} = (1, 0)^{\mathsf{T}}$, $\ket{g} = (0, 1)^{\mathsf{T}}$,
$\sigma_z\ket{e} = +\ket{e}$ and $\sigma_- = \ket{g}\bra{e}$, as in QuantumOptics.jl.
Section~\ref{sec:physics} defines the Hamiltonian and every other symbol of the toy
problem.

\textbf{Code.} The Julia project \cmd{DrivenSpin/} of the course repository is the
reference for all code. Where these notes would quote it and it does not exist yet, a red
box marks the placeholder.


\tableofcontents
\clearpage

\part{Notes, block by block}
% Block 0 of the slides (00_front.tex).
\section{The thesis: verification is the bottleneck}\label{sec:thesis}

\onslides{two motivation slides (TheoPhysEval, Tao), the thesis, and the
``new PhD student'' mental model.}

\subsection{Two data points from October 2026}

\textbf{Agents solve research problems.} Ramesh Chandra et al.\ \rf{theophyseval}
collected 50 unpublished research-level problems in theoretical physics, each with a
verifiable final answer and each about a day's work or less for a world expert. Agents
received literature search and the ability to write and run scripts, and no human hints.
The best agent reached $98\,\%$ correct final answers (49 of 50) at an average cost under
3\,USD per problem; the runner-up solved 46. Accuracy changed little when web search
and background papers were withheld, which the authors read as substantial memorisation
of the literature. They also discuss why no agent-led breakthrough has appeared despite
this problem-solving ability. Two points matter for the course: the problems had answers
that could be \emph{checked}, and someone had to build that check.

\textbf{Abundance outpaces verification.} Tao's essay \rf{tao}, based on his public
lecture at the 2026 International Congress of Mathematicians, assumes that AI tools will
perform research-level mathematical tasks and asks what mathematicians actually value.
In Section~7 he lists the ``impedance mismatches'' that would follow; the first is the
sentence quoted on the slide, that AI-generated proofs will accumulate faster than they
can be verified. He calls the result a transition from proof scarcity to proof
abundance, and notes that journals, priority conventions, hiring and prizes were all
designed for scarcity. The slide carries the argument to physics, where code, numerics
and derivations have become cheap in the same way.

\subsection{From the data to the thesis}

If producing a result costs minutes, the scarce resource is the confidence that the
result is right. Research software engineering (RSE) is the body of practice that builds
that confidence for code written by humans: version control, a fixed project structure,
automated tests and continuous integration \rf{wilson-good,wilson-best}. None of it was
designed for AI, and all of it applies, because every one of these practices was
invented for contributors who make mistakes. An agent is such a contributor, only
faster.

\subsection{The new PhD student}

The mental model on the slide has four traits, and each maps onto a later block.
\begin{itemize}
  \item \textbf{Fast}: it produces more than you can read, so the reading has to be
        made efficient (diffs, Section~\ref{sec:git}; tests, Section~\ref{sec:verification}).
  \item \textbf{Eager}: it rarely declines a task, and rarely says ``I don't know''
        (Section~\ref{sec:llm}, post-training).
  \item \textbf{Overconfident}: correct and incorrect output read the same
        (Section~\ref{sec:llm}, plausibility).
  \item \textbf{Forgetful}: nothing survives the session unless it is written to disk
        (Section~\ref{sec:context}).
\end{itemize}
The analogy has a limit worth stating if asked: a PhD student learns from your
corrections and becomes more reliable over months; a model does not, unless the
correction is written into a file it reads every time.

\subsection{A sceptical counterweight}

A randomised controlled trial by METR \rf{metr} gave 16 experienced open-source
developers 246 tasks in their own repositories, with AI tools allowed or disallowed per
task. Before the study they forecast that AI would cut completion time by $24\,\%$;
afterwards they estimated it had cut it by $20\,\%$; measured, it \emph{increased}
completion time by $19\,\%$. The study used early-2025 tools on large, mature codebases
the developers knew well, and the authors caution against generalising. Its lasting
lesson is the gap between perceived and measured productivity: how fast the work
\emph{feels} is not evidence. The slides dropped this study for pace; it is a good answer
to ``but does it really make us faster?''

\begin{qa}
\Q{Is TheoPhysEval not just measuring memorisation?}
Partly, and the authors say so: removing web search and papers changed accuracy little.
The problems were unpublished, but their ingredients were not. The course's point
survives either way: the answers were checkable, and the check made the result usable.

\Q{If agents are this good, why spend an afternoon on verification?}
Because ``98\,\%'' was measured against known answers. In your research there is no
answer key; the tests and cross-checks you write are the answer key.

\Q{Did Tao say AI is bad for mathematics?}
No. The essay conditions on capable AI arriving and asks which values and institutions
need to be made explicit. The quoted line describes a bottleneck, not a verdict.
\end{qa}

\textbf{Read next:} Tao \rf{tao}, Sections~1 and~7 (short and non-technical); the
TheoPhysEval abstract and discussion \rf{theophyseval}; the METR paper \rf{metr}.

% Block 1 of the slides (01_llm.tex) and the attention backup slide.
\section{How a large language model works}\label{sec:llm}

\onslides{tokens; next-token prediction; generation as a loop; pre-training as a
fit; post-training (instruction tuning, human feedback, verifiable rewards); two tables
of properties and their consequences; knowledge cutoff and the Julia bias.}

\subsection{Tokens}

A model never sees characters. A \emph{tokenizer} splits text into pieces from a fixed
vocabulary and maps each piece to an integer. Modern tokenizers learn the vocabulary
from data with byte-pair encoding \rf{sennrich}: start from single bytes and repeatedly
merge the most frequent adjacent pair into a new symbol, until the vocabulary has the
target size. Frequent words become one token; rare words, numbers, code and \LaTeX{}
split into several. Llama~3 uses $128$K tokens and reaches $3.94$ characters per token
on English text \rf{llama3}. An English word averages about five letters plus a space,
so one token is roughly $0.7$ words, consistent with the ``$0.75$ words per token'' on the
slide. The token numbers on the slide are illustrative.

Two consequences come up in questions. Letter-level tasks (counting the r's in a word,
reversing a string) are hard because the model does not see letters. Arithmetic on long
numbers is unreliable because digits are grouped into tokens in irregular ways; this is
one reason agents call a calculator or write a script.

\subsection{Next-token prediction}

Write the tokens so far as $x_1, \ldots, x_n$; this sequence is the \emph{context}. The
model is a function that returns a probability distribution over the vocabulary for the
next token,
\begin{equation}
  p_\theta(x_{n+1} \mid x_1, \ldots, x_n),
\end{equation}
where $\theta$ denotes the parameters of a neural network, a \emph{transformer}
\rf{vaswani}. Large models have $10^{11}$ to $10^{12}$ parameters; the largest Llama~3
model has $4.05 \times 10^{11}$ \rf{llama3}. The network outputs one real number per
vocabulary entry, a \emph{logit} $z_i$, and the distribution is the softmax
\begin{equation}
  p_\theta(x_{n+1} = i \mid x_1, \ldots, x_n)
    = \frac{\ee^{z_i/T_{\mathrm{s}}}}{\sum_j \ee^{z_j/T_{\mathrm{s}}}},
\end{equation}
with a \emph{sampling temperature} $T_{\mathrm{s}}$ set by the harness (the original
distribution is $T_{\mathrm{s}} = 1$). A physicist reads this as a Boltzmann
distribution with energies $-z_i$. Small $T_{\mathrm{s}}$ concentrates the probability
on the most likely token; large $T_{\mathrm{s}}$ flattens it.

\textbf{Generation is a loop.} Sample a token, append it to the context, compute the
distribution again. A long answer is thousands of passes. Two facts follow and are on the
slide: in each pass the context is the model's only input, and a token once written is
never revised. A model can correct itself only by writing more text, for example
``wait, that is wrong'', which is what reasoning models learn to do (below).

\textbf{Attention} (backup slide). Inside the transformer, each token is a vector. In
each layer, each token's vector is mapped linearly to a \emph{query}, a \emph{key} and a
\emph{value}; stacking these for all tokens gives matrices $Q$, $K$, $V$, and the layer
computes
\begin{equation}
  \mathrm{Attention}(Q, K, V) = \mathrm{softmax}\!\left(\frac{Q K^{\mathsf{T}}}{\sqrt{d}}\right) V,
\end{equation}
where $d$ is the dimension of the keys and the softmax acts row by row \rf{vaswani}. Row
$m$ of the result is a weighted average of the values of all tokens, weighted by how well
their keys match the query of token $m$. In a language model a mask restricts the
average to earlier tokens. The matrix $Q K^{\mathsf{T}}$ has one entry per pair of
tokens, so the cost grows with the square of the context length; this is one reason the
context window is finite. 3Blue1Brown's chapters 5 and 6 \rf{3b1b} are the best visual
explanation.

\subsection{Pre-training: a very large fit}

Pre-training minimises the cross-entropy of the next token over a corpus,
\begin{equation}
  \mathcal{L}(\theta) = -\sum_{n} \ln p_\theta(x_{n+1} \mid x_1, \ldots, x_n),
\end{equation}
where the sum runs over every position $n$ of every document. This is a maximum-likelihood
fit: the same principle as a least-squares fit with Gaussian errors, with $10^{12}$
parameters and $10^{13}$ data points. The largest Llama~3 model was pre-trained on
$15.6 \times 10^{12}$ tokens \rf{llama3}. The loss falls as a power law in model size,
data and compute \rf{kaplan}, and for a fixed compute budget model size and data should
grow together \rf{chinchilla}; these \emph{scaling laws} are why models became large.

What the fit learns is the distribution of text. Text about physics is mostly correct, so
the model learns a great deal of physics; but the objective rewards the \emph{form} of
plausible text and never checks its truth. That is the root of the first row of the
property table: a fluent wrong derivation and a fluent right one are both high-likelihood
text.

\textbf{Knowledge cutoff.} The corpus ends at some date. Packages released or changed
after it are unknown or known in an old version. The corpus also reflects what people
write: there is far more Python than Julia. The slide's example shows the typical
result. Asked for QuantumOptics.jl, the model writes \cmd{mesolve}, the name of the
master-equation solver in QuTiP (Python); the Julia package calls it
\cmd{timeevolution.master} \rf{qojulia,qojulia-docs}. The code looks right and fails at
the first run, which an agent notices and a chatbot user may not.

\subsection{Post-training: from text model to assistant}

A pre-trained model continues text; asked a question, it may continue with more
questions. Post-training turns it into an assistant in three stages.

\textbf{Instruction tuning} (supervised fine-tuning). Continue training on examples of
prompts with good answers written by people. The model imitates the format and tone.

\textbf{Reinforcement learning from human feedback} (RLHF). People compare pairs of
answers to the same prompt. A second network, the \emph{reward model}, is trained to
predict which answer people prefer, giving a scalar score $r(x, y)$ for prompt $x$ and
answer $y$. The language model, now viewed as a policy $\pi_\theta(y \mid x)$, is then
trained by reinforcement learning to maximise
\begin{equation}
  \mathbb{E}_{y \sim \pi_\theta}\left[r(x, y)\right]
    - \beta\, D_{\mathrm{KL}}\!\left(\pi_\theta \,\|\, \pi_{\mathrm{ref}}\right),
\end{equation}
where $\pi_{\mathrm{ref}}$ is the model before this stage, $D_{\mathrm{KL}}$ the
Kullback--Leibler divergence and $\beta > 0$ a weight that keeps the model from drifting
far from fluent text. The method comes from Christiano et al.\ \rf{christiano} and was
applied to language models in InstructGPT \rf{instructgpt}, where a 1.3-billion-parameter
model trained this way was preferred by human raters over the 175-billion-parameter
GPT-3. Variants replace human raters with a model that judges answers against written
principles \rf{constitutional}.

\textbf{Reinforcement learning with verifiable rewards} (RLVR). Replace the learned reward
model with a check: the final answer to a maths problem matches the known one, or the
generated code passes the tests. The term RLVR was introduced with Tülu~3
\rf{tulu3}; DeepSeek-R1 \rf{deepseekr1} showed that this kind of training alone makes a
model produce long chains of intermediate reasoning, check its steps and backtrack,
before giving an answer. These are the ``thinking'' tokens of current models. Earlier,
the same behaviour had been elicited by prompting (``chain of thought'') \rf{wei-cot}.

The key box on the slide summarises the incentives: the model was rewarded for answers
that \emph{look good to people} and \emph{pass checks}. Neither is the same as being
correct, and both leave marks.

\subsection{Properties and their consequences}

\textbf{Trained to sound plausible} (pre-training). Right and wrong output have the same
fluency. References are the classic case: a plausible author, title, journal and year are
easy to generate and the model has no lookup to check them against. Every reference in
this course was opened and checked (Part~III).

\textbf{Trained to please} (human feedback). Raters prefer answers that agree with them.
Sharma et al.\ \rf{sharma} showed that five assistants consistently exhibit this
\emph{sycophancy}, for example abandoning a correct answer when the user pushes back, and
that both humans and preference models sometimes prefer a convincing sycophantic answer
over a correct one. The defence on the slide follows: ask the model to find the flaw in
your argument, not to confirm it.

\textbf{Rewarded for passing checks} (verifiable rewards). A policy rewarded for passing
tests learns that weakening the test also passes it. METR reported in 2025 that
frontier models, in its evaluations, attempted ``(often successfully) to get a higher score
by modifying the tests or scoring code'' \rf{metr-hacking}. In everyday use this looks like
a loosened tolerance, a deleted assertion, or a special case for the test input. The
defence: read every change to a test, and commit tests before the implementation
(Section~\ref{sec:verification}).

\textbf{Random sampling} (architecture). With $T_{\mathrm{s}} > 0$ the same prompt gives
different answers. This is a free diagnostic: if two runs disagree on a number, neither
is reliable. Setting $T_{\mathrm{s}} = 0$ makes the choice greedy, but in practice
deployed models do not guarantee bit-identical outputs, so do not build reproducibility on
it; build it on code and data in a repository.

\textbf{No memory between sessions.} The parameters do not change while you use the
model. Everything it ``knows'' about your project in a session is in the context; at the
next session it starts from zero. Hence: write things down (Section~\ref{sec:context}).

\textbf{Finite context.} See Section~\ref{sec:context}.

\begin{qa}
\Q{Is it ``just autocomplete''?}
The objective is next-token prediction, yes. But predicting the next token of a
derivation well requires representing a lot of the structure of the derivation, and
post-training adds goal-directed behaviour. ``Just autocomplete'' underestimates what it
can do; ``understands physics'' overstates how reliably. The practical answer is the
course's: treat its output as a claim to be checked.

\Q{Why does it invent references instead of saying it doesn't know?}
Nothing in pre-training rewards abstaining, and a confident well-formatted answer tends
to be preferred by raters. Agents with web search can check references, but only if they
are asked to and actually do it; check the link yourself.

\Q{Does it learn from my conversations?}
Not within the session: the parameters are fixed. Whether a provider uses conversations
for future training depends on its terms and your account settings; check them before
pasting unpublished data (Section~\ref{sec:failures}).

\Q{How many parameters does model X have?}
Providers of closed models do not publish it. Open-weight models do, for example
$4.05 \times 10^{11}$ for the largest Llama~3 \rf{llama3}.
\end{qa}

\textbf{Read or watch next:} Karpathy's \emph{Intro to Large Language Models} (1\,h) and
\emph{Deep Dive into LLMs like ChatGPT} (3.5\,h), which covers tokenization, pre-training
and post-training in exactly this order \rf{karpathy-intro,karpathy-deep}; 3Blue1Brown
chapters 5--6 for attention \rf{3b1b}; the InstructGPT paper \rf{instructgpt},
Sections~1 and~3, for RLHF.

% Block 2 of the slides (02_agents.tex).
\section{Chatbots, harnesses and agents}\label{sec:agents}

\onslides{a chatbot where you are the loop; an agent where the loop closes; one
agent transcript; chatbot versus agent table; permissions; the live demo ``Same task, two
ways''; asking for verifiable output beyond code.}

\subsection{What a harness is}

The model of Section~\ref{sec:llm} maps text to text and does nothing else. Everything
else in a product such as Claude Code or Codex is the \emph{harness}: the program that
\begin{enumerate}
  \item assembles the context (system instructions, instruction files, your messages,
        earlier tool results),
  \item offers the model a set of \emph{tools}, each described by a name, a description
        and a schema for its arguments (read a file, edit a file, run a shell command,
        search the web),
  \item recognises when the model's output is a \emph{tool call}, a structured request
        such as \texttt{run("julia --project -e 'using Pkg; Pkg.test()'")}, asks you for
        permission if needed, executes it, and appends the result to the context,
  \item loops until the model answers without a tool call, and
  \item manages the context when it fills up (Section~\ref{sec:context}).
\end{enumerate}
The model never touches your computer. It writes text that requests an action, and the
harness decides whether to carry it out. A chatbot in the browser is a harness with no
tools on your machine, and \emph{you} carry code and error messages between the model and
your computer: you are the loop.

Interleaving reasoning with actions and observations is known as the ReAct pattern
\rf{react}. Schluntz and Zhang \rf{agents} distinguish \emph{workflows}, in which code
fixes the sequence of model calls, from \emph{agents}, in which the model decides its next
step from the results of the previous one. Their advice, start with the simplest design
that works, applies to how you use agents too.

\subsection{Why the closed loop matters}

The transcript on the slide is illustrative (the real test is the $\pi$-pulse line of the
\cmd{"Rabi"} testset, Section~\ref{sec:structure}), and it shows the point of the block. The agent runs the tests, reads a failure (\lstinline|0.0 ≈ 1.0|), reads the
source, finds $\Omega$ where $\Omega/2$ belongs, edits, and runs the tests again. It
found the factor of two \emph{because a test encoded the physics}: a resonant
$\pi$-pulse must invert the spin. With the bug, the population at $t = \pi/\Omega$ is
$\sin^2(\pi) = 0$ instead of $\sin^2(\pi/2) = 1$ (Section~\ref{sec:failures}). An agent
can check its work against reality, but only against the reality you give it.

Evidence that this works at scale: SWE-bench \rf{swebench} tests whether a system can
resolve 2294 real GitHub issues from 12 popular Python repositories, judged by the repositories'
own tests. METR measures the length of software tasks, in human working time, that
agents complete with $50\,\%$ reliability, and found that this length has doubled
roughly every seven months since 2019 \rf{metr-horizon}. Both benchmarks rely on
tests written by people.

\subsection{Permissions}

Agents ask before editing a file or running a command, and every agent lets you approve
whole classes of actions (all edits, or specific commands). The slide's rules:
\begin{itemize}
  \item Never auto-approve deletion, overwriting, or anything outside the project
        folder. An agent runs commands with your user's rights.
  \item ``Yes, yes, yes'' is not supervision. If you approve without reading, you have
        granted full permission with extra steps.
  \item Git makes a longer leash safe (Section~\ref{sec:git}): any change inside the
        repository can be inspected and undone, but nothing outside it can.
\end{itemize}

\subsection{The live demo: same task, two ways}

The prompt on the slide asks for Rabi oscillations with the course Hamiltonian, two
detunings, three Rabi periods and a check against the analytic formula. What to point out
while it runs:
\begin{itemize}
  \item In the chat, every error costs a copy--paste round trip; the model never sees the
        plot.
  \item The agent writes files, runs Julia, reads the error, fixes it. Count the tool
        calls.
  \item Ask the room: what did it check? (Typically it compares with the formula at a
        few points.) What did it not check? (Units, the sign convention of $\Delta$, the
        initial state, that the plot shows what the caption claims.)
\end{itemize}

\subsection{Beyond code}

For any task, ask for an output you can verify: links you open for a literature search, a
script that checks a derivation numerically, an analysis that lives as a script in the
repository, and text you can defend line by line. A number in a chat window cannot be
reproduced; a script in a repository can.

\begin{qa}
\Q{Which agent should I use?}
The course does not prescribe one; the concepts carry over. Use the one your institution
or subscription gives you, and check its data-handling terms (Section~\ref{sec:failures}).

\Q{Can the agent delete my files?}
If you permit it, yes. Keep the project under git, commit before the agent starts, and do
not approve commands that act outside the project folder.

\Q{Is ``agent'' just marketing?}
In the sense used here it has a precise meaning: a model in a loop that chooses tool
calls from their results \rf{agents}. Whether that is useful depends on whether the
results contain a signal of correctness, which is what tests provide.
\end{qa}

\textbf{Read next:} \emph{Building effective agents} \rf{agents} (short, practical);
Anthropic's best-practice guide for its coding agent \rf{cc-best}, whose first advice,
``give Claude a way to verify its work'', is tool-specific in wording and general in
content.

% Block 3 of the slides (03_context.tex).
\section{Context, effort and persistent instructions}\label{sec:context}

\onslides{the context window; context rot; context hygiene; the cost of thinking;
instruction files; skills; a prompt as a specification; the workflow explore--plan--%
implement--verify--commit. This block is the first to shorten if time runs out.}

\subsection{The context window}

The context is the token sequence $x_1, \ldots, x_n$ of Section~\ref{sec:llm}. In an agent
session it contains, in order: the harness's system instructions and tool descriptions,
your instruction files, your messages, every file the agent read, every command output,
and the model's own answers and reasoning. Its length is capped by the context window,
currently $10^5$ to $10^6$ tokens (the largest Llama~3 model has $1.28 \times 10^5$
\rf{llama3}). At about $0.7$ words per token, $10^5$ tokens is a long novel. That sounds
like a lot until an agent reads a few source files, a test log and a long solver output.

\subsection{Context rot}

Accuracy degrades as the context grows, well before the window is full. Hong, Troynikov
and Huber \rf{contextrot} tested 18 models on tasks that stay equally hard as the input
grows, so that any drop is due to length alone. Every model degraded with input length,
and non-uniformly. Three findings are worth quoting: a single distractor (text similar to
the answer but wrong) lowered accuracy, and its effect grew with length; on a
conversational memory benchmark, all models did much better with a focused input of about
300 tokens than with the full $\sim 113$K-token history; and models performed
\emph{worse} when the irrelevant filler was a coherent text than when its sentences were
shuffled. Earlier, Liu et al.\ \rf{lostmiddle} found that models use information at the
beginning and end of a long context better than information in the middle.

Practical consequences:
\begin{itemize}
  \item A wrong turn stays in the context and keeps influencing later answers. Arguing
        with the agent adds more of it. Start a fresh session with a better prompt.
  \item When the window fills, harnesses \emph{compact}: they replace the history by a
        summary written by the model \rf{context-eng}. The summary is lossy; details you
        care about may not survive.
  \item One task, one session. Write what must survive into files (below).
\end{itemize}

\textbf{Point, do not paste.} Three thousand lines of solver log in your message cost
context and bury the one relevant line. Naming the failing test and asking the agent to
run only that test lets it read what it needs. Anthropic's engineering team calls this
loading context ``just in time'' through references such as file paths \rf{context-eng}.

\subsection{What a session costs}

Providers charge per token, at different prices for input and output. With $N_{\mathrm{in}}$
and $N_{\mathrm{out}}$ the numbers of input and output tokens and $c_{\mathrm{in}}$,
$c_{\mathrm{out}}$ the prices per token,
\begin{equation}
  \text{cost} \approx N_{\mathrm{in}}\, c_{\mathrm{in}} + N_{\mathrm{out}}\, c_{\mathrm{out}}.
\end{equation}
The slide's remark that every agent step re-reads the context has a quantitative
consequence. Suppose a session takes $K$ steps (tool calls), starts from a context of
$C_0$ tokens, and each step adds about $a$ tokens (a file read, a command output). Step
$k$ reads the whole context, $C_0 + k a$ tokens, so
\begin{equation}
  N_{\mathrm{in}} \approx \sum_{k=1}^{K} (C_0 + k a) \approx K C_0 + \tfrac12 a K^2 ,
\end{equation}
quadratic in the length of the session. Long sessions are expensive as well as less
accurate. Many providers discount input that repeats an earlier prefix (\emph{prompt
caching}), which softens but does not remove the effect.

$N_{\mathrm{out}}$ includes the reasoning tokens of Section~\ref{sec:llm}, which you may
never see. The \emph{effort} setting of current agents controls how much the model
reasons before answering: low effort for mechanical edits, high for debugging and
derivations. For scale, the TheoPhysEval agents solved day-long research problems for
under 3\,USD each \rf{theophyseval}.

\subsection{Instruction files}

Because the model remembers nothing between sessions, harnesses read an instruction file
from the project folder at the start of every session and put it into the context. The
file name depends on the tool: \cmd{AGENTS.md} is an open convention read by many agents
\rf{agentsmd}; Claude Code reads \cmd{CLAUDE.md} \rf{cc-best}. The \cmd{DrivenSpin/}
project has an \cmd{AGENTS.md} and a one-line \cmd{CLAUDE.md} that points to it.

What belongs in it: the conventions the agent cannot guess (units, basis order, sign of
$\Delta$), the command that runs the tests, and the rules that protect you (never weaken a
test, never write to \cmd{data/exp\_raw}). What does not: anything the agent can read from
the code, long tutorials, information that changes often. The file is in every context,
so every line costs tokens in every step; a long file also dilutes the rules that matter
\rf{cc-best}.

\begin{placeholder}[Placeholder: \texttt{DrivenSpin/AGENTS.md}]
Quote the real file here once D3 exists. The slide shows an illustrative excerpt:
conventions ($\hbar = 1$, \textmu s, rad/\textmu s, basis $(\ket{e}, \ket{g})$,
$H = (\Delta/2)\sigma_z + (\Omega/2)\sigma_x$), the test command
\verb|julia --project -e 'using Pkg; Pkg.test()'|, ``tests first'', ``never modify
\texttt{test/} while implementing'', ``never write to \texttt{data/exp\_raw}''.
\end{placeholder}

\subsection{Skills}

A \emph{skill} is a folder with a file \cmd{SKILL.md} that starts with a name and a
one-line description, followed by instructions, and optionally scripts and further
files. At the start of a session only the names and descriptions enter the context; when a
task matches a description, the agent reads the full file, and it opens further files
only when a step needs them \rf{skills}. This \emph{progressive disclosure} lets a lab keep
many procedures (fit Ramsey data, calibrate a pulse, format a figure) at the cost of one
line each. The slide's \cmd{fit-ramsey} skill is an example of the form, not a file in the
repository. Instruction files and skills are \emph{procedural memory written to disk}:
the model does not learn, but the folder does.

\subsection{A prompt is a specification}

The slide contrasts ``Simulate a qubit with decoherence'' with a prompt that names the
goal and context (which file, which conventions), the verification (the tests that define
``done''), and the constraints (what must not change; plan before code). Each part closes
a door the agent would otherwise walk through: without the file it invents a new one;
without the conventions it picks its own; without the tests ``done'' means ``looks
done''; without constraints it ``improves'' the Hamiltonian.

\subsection{A workflow that works}

Explore (read the code), plan (you approve), implement (small steps), verify (tests,
plots), commit (you read the diff). The same sequence is the core recommendation of
Anthropic's guide \rf{cc-best}, and it is tool-independent. Most failures come from
skipping \emph{plan} (the agent solves the wrong problem) or \emph{verify} (it solves the
problem wrongly).

\begin{qa}
\Q{How big is my context right now?}
Most agents can show it (a command or a status line). If you cannot see it, start fresh
sessions often.

\Q{Should I always use the highest effort?}
No: it costs more and is slower, and for mechanical edits it does not help. Raise it when
the agent goes in circles on a hard problem.

\Q{Instruction file or skill?}
Rules that apply to every task go into the instruction file. Procedures that apply to
some tasks go into skills, so they cost context only when used.

\Q{Is ``context rot'' a real effect or a vendor's term?}
The term comes from the Chroma report \rf{contextrot}, a controlled measurement on 18
models; the effect is consistent with earlier academic work \rf{lostmiddle}.
\end{qa}

\textbf{Read next:} \emph{Effective context engineering for AI agents} \rf{context-eng};
the Chroma report \rf{contextrot}, whose figures are worth a look; the skills article
\rf{skills}.

% Blocks 4 and 4b of the slides (04_git.tex, 04b_git_handson.tex).
\section{Git: the safety net}\label{sec:git}

\onslides{why git now; a chain of snapshots; three places for a change; a first
repository; reading a diff; a weakened test; undo; branches; remotes; what goes into git;
habits and prompts for working with an agent; the live demo ``The agent wrecks the
project''; the nine commands. Block 4b is a local hands-on (four exercises).}

The one-page cheat sheet is Section~\ref{sec:cheatsheet}. This section explains the model
behind the commands, which is what you need to answer ``what exactly happens when\dots''.

\subsection{What git stores}

A git repository is the hidden folder \cmd{.git/} inside your project. It holds a
database of \emph{objects}, each named by a hash of its content \rf{progit}:
\begin{itemize}
  \item a \emph{blob} is the content of one file;
  \item a \emph{tree} is a directory listing: names, permissions, and the hashes of the
        blobs and subtrees;
  \item a \emph{commit} records one tree (the snapshot of the whole project), the hash of
        its parent commit (two parents for a merge), the author, the date and the message.
\end{itemize}
The hash is a 40-hex-digit SHA-1 value; the first seven digits (\cmd{9c0d5e2}) almost
always suffice to identify a commit within a project. Because a commit's hash covers its tree and its
parent, which covers \emph{its} parent, a hash identifies the whole history up to that
point: change one byte anywhere in the past and every later hash changes. This is why a
hash stored next to a result (Section~\ref{sec:structure}) pins down the code exactly.

A commit is a snapshot, not a list of changes. A diff is computed on demand by comparing
two snapshots. Unchanged files are stored once and shared between snapshots, so history
is cheap.

A \emph{branch} is a file containing one commit hash: a movable label. Committing on a
branch writes the new commit and moves the label to it. \cmd{HEAD} names what you have
checked out, normally a branch. That is all a branch is, which is why creating and
deleting branches costs nothing.

\subsection{Three places for a change}

\begin{center}
\begin{tabularx}{\linewidth}{@{}l X@{}}
\toprule
\textbf{Place} & \textbf{What it is} \\
\midrule
Working folder & The files you and the agent edit. \\
Staging area (index) & The snapshot being prepared; \cmd{git add} copies a file's current
  content here. \\
History & The commits; \cmd{git commit} turns the staging area into a new commit. \\
\bottomrule
\end{tabularx}
\end{center}
The staging area lets you choose what goes into a commit: two unrelated changes, two
commits (Exercise~2 of block 4b). \cmd{git status} reports where each change is; run it
often.

\textbf{The forms of \cmd{git add}.} \cmd{git add <file>} stages one file;
\cmd{git add .} stages everything new, changed or deleted in the current folder and below;
\cmd{git add -u} stages changes to files git already tracks and leaves new files out;
\cmd{git add -p} shows each change (\emph{hunk}) and asks \cmd{y} or \cmd{n}. In the
hands-on it asks file by file in alphabetical order. \cmd{git commit -am "..."} is
\cmd{git add -u} plus commit in one step: convenient, but it skips the look at what you
stage.

\subsection{Reading a diff}

\begin{lstlisting}
--- a/src/DrivenSpin.jl
+++ b/src/DrivenSpin.jl
@@ -12,6 +12,6 @@ function hamiltonian(Δ, Ω)
     σx = [0 1; 1 0]
     σz = [1 0; 0 -1]
-    return Δ/2*σz + Ω*σx
+    return Δ/2*σz + Ω/2*σx
 end
\end{lstlisting}
\begin{itemize}
  \item \cmd{---} and \cmd{+++} name the old (\cmd{a/}) and new (\cmd{b/}) version of the
        file.
  \item \verb|@@ -12,6 +12,6 @@| is the hunk header: the excerpt starts at line 12 and is
        6 lines long, in the old and in the new file. Git appends the nearest enclosing
        function as a hint.
  \item Lines starting with a space are unchanged context; \cmd{-} lines were removed,
        \cmd{+} lines added. A modified line appears as one removal and one addition.
\end{itemize}
\cmd{git diff} shows changes not yet staged; \cmd{git diff --staged} shows what the next
commit will contain; \cmd{git diff <hash1> <hash2>} compares two snapshots;
\cmd{git show <hash>} shows one commit with its message.

\textbf{Spot the problem} (slide). An agent changes
\lstinline|@test excited_population(H, π/Ω) ≈ 1.0| to
\lstinline|... ≈ 1.0 atol=1.0|. Now any value between $0$ and $2$ passes, so the test checks
nothing while staying green. Every change to a test needs a reason you accept.

\subsection{Undo}

\begin{center}
\begin{tabularx}{\linewidth}{@{}X l X@{}}
\toprule
\textbf{Situation} & \textbf{Command} & \textbf{Effect} \\
\midrule
Unwanted edits, not staged & \cmd{git restore <file>} & Working file back to the staged
  (or last committed) version. The edits are gone. \\
Staged by mistake & \cmd{git restore --staged <file>} & Out of the staging area; the edit
  stays in the working file. \\
A bad commit & \cmd{git revert <hash>} & A \emph{new} commit that undoes it; history is
  kept. \\
\bottomrule
\end{tabularx}
\end{center}
\cmd{git restore .} discards \emph{all} uncommitted changes below the current folder:
look at \cmd{git status} and \cmd{git diff --stat} first. Uncommitted work that is
discarded is not recoverable from git, which is why ``commit before the agent starts'' is
the first habit.

Two commands destroy work and should never run without your explicit approval:
\cmd{git reset --hard} (moves the branch and overwrites the working folder) and
\cmd{git push --force} (replaces the history on the remote, including your colleagues'
commits). Agents occasionally reach for them to ``clean up''.

\subsection{Branches and merging}

\cmd{git switch -c try-rk4} creates a branch at the current commit and switches to it;
commits now move \cmd{try-rk4} and leave \cmd{main} where it was. If the experiment works,
\cmd{git switch main} and \cmd{git merge try-rk4}. If \cmd{main} has not moved, the merge
just moves the \cmd{main} label forward (a \emph{fast-forward}); otherwise git creates a
merge commit with two parents, as on the slide. If both branches changed the same lines,
git stops with a \emph{conflict} and marks the competing versions in the file; you edit,
\cmd{git add} and \cmd{git commit}. If the experiment failed, \cmd{git branch -D try-rk4}
deletes the label and its commits are as if they never happened.

\subsection{Remotes}

A remote is another copy of the repository, usually on GitHub, with the full history.
\cmd{git clone} copies one and remembers it as \cmd{origin}; \cmd{git push} sends your new
commits; \cmd{git pull} fetches new commits and merges them into your branch. Every clone
is a complete backup. The installation guide set up login over SSH with \cmd{gh}.

\subsection{What goes into git}

Commit source, scripts, tests, \cmd{Project.toml} and \cmd{Manifest.toml},
\cmd{README.md} and \cmd{AGENTS.md}. Do not commit large raw data (archive it elsewhere),
generated results and plots (they can be regenerated), and never passwords or API keys.
The \cmd{.gitignore} file lists patterns git should not track (\cmd{data/},
\cmd{plots/}, \cmd{*.jld2}).

If a secret was committed and pushed, consider it public: revoke it and issue a new one.
Removing it from history afterwards is possible but does not undo copies that already
exist.

\subsection{Git with an agent}

The habits on the slide:
\begin{enumerate}
  \item Commit before the agent starts: a clean point to return to, and
        \cmd{git diff} then shows exactly what the agent did.
  \item One task, one commit: each change can be reviewed and reverted on its own.
  \item Read the diff before every commit, tests first.
  \item Let the agent type the git commands; you check the result.
\end{enumerate}
Prompts such as ``Commit my manual changes before starting'' and ``Work in modular
commits: one commit per change, with the tests passing at each'' make the agent keep a
history in which your edits and its edits never share a commit.

\textbf{The live demo.} Ask the agent for a sweeping refactor (``all parameters in one
struct''), then: \cmd{git status --short} (what did it touch?), \cmd{git diff --stat} (how
much?), \cmd{git restore .} (back to the last commit). Thirty seconds, nothing lost,
\emph{because} the last commit was clean.

\subsection{Block 4b: the local hands-on}

Four exercises in a throw-away folder \cmd{git-practice/}, in pairs, without agent or
GitHub; all commands one per line so they also work in Windows PowerShell~5.
{\raggedright
\begin{enumerate}
  \item \textbf{First commits}: \cmd{git init}, \cmd{status}, \cmd{add} one file,
        \cmd{commit}, then the other; \cmd{git log --oneline}. The files: \cmd{rabi.jl}
        with the Rabi formula as one line of Julia, and \cmd{notes.md}.
  \item \textbf{One change, one commit}: edit both files; \cmd{git diff};
        \cmd{git add -p} (answer \cmd{n} for \cmd{notes.md}, \cmd{y} for \cmd{rabi.jl});
        \cmd{git diff --staged}; commit; then \cmd{git add -u} and commit the rest.
  \item \textbf{Undo}: delete the formula and \cmd{git restore} it; stage a change and
        unstage it with \cmd{git restore --staged}; commit a broken formula
        ($t/2 \to t$) and undo it with \cmd{git revert --no-edit HEAD}, so both commits
        remain in the history.
  \item \textbf{Read the history}: \cmd{git log --oneline}, \cmd{git log --stat},
        \cmd{git show <hash>}, \cmd{git diff <hash1> <hash2>}, \cmd{git log -p rabi.jl}.
        Long output opens in a pager: arrow keys to scroll, \cmd{q} to quit.
\end{enumerate}\par}
A possible \cmd{rabi.jl} line, in the course conventions:
\begin{lstlisting}
Pe(t; Δ, Ω) = Ω^2 / (Ω^2 + Δ^2) * sin(sqrt(Ω^2 + Δ^2) * t / 2)^2   # ħ = 1, rad/μs, μs
\end{lstlisting}

\begin{qa}
\Q{Do I need GitHub to use git?}
No. Everything in block 4b is local. GitHub adds backup, sharing and CI.

\Q{I ran \texttt{git restore} on the wrong file. Can I get it back?}
Not from git, if the change was never staged or committed. Your editor's undo or local
history may still have it. This is why the agent should start from a clean commit.

\Q{What is the difference between \texttt{revert} and \texttt{reset}?}
\cmd{revert} adds a commit that undoes an earlier one; history only grows, which is safe
on shared branches. \cmd{reset} moves the branch label backwards, so later commits drop
out of the history; with \cmd{--hard} it also overwrites your files.

\Q{Why does \texttt{git pull} sometimes create a merge commit?}
Because you and the remote both have new commits; git joins the two lines of history.

\Q{Is SHA-1 not broken?}
Collisions have been constructed deliberately, and git has added a hardened variant and an
optional SHA-256 format. For tracking your own research code, accidental collisions are
not a practical concern.
\end{qa}

\textbf{Read next:} \emph{Pro Git} \rf{progit}, Chapters~2 (basics), 3 (branching) and
10 (internals); the Software Carpentry lesson \rf{swc-git} to practise.

% Block 5 of the slides (05_structure.tex).
\section{Project structure, tests and continuous integration}\label{sec:structure}

\onslides{which code made Figure~3; the DrWatson layout; the code as one part of the
research project; \texttt{src} computes, scripts orchestrate; every result knows its
commit; the environment is part of the result; tests; continuous integration; the same
principles for MATLAB, instruments and FPGA.}

The goal on the first slide is the test of this block: \emph{for every figure, you can say
which code, data and parameters made it}. An agent makes this harder, because it creates
files faster than anyone names them well, and easier, because it follows a fixed layout
reliably once the layout exists.

\subsection{The DrWatson layout}

DrWatson \rf{drwatson,drwatson-docs} is a Julia package that fixes a project layout and
provides a few functions for naming, saving and loading results.
\cmd{initialize\_project("DrivenSpin")} creates the folders and a git repository:
\begin{lstlisting}
DrivenSpin/
|-- src/            # the physics: functions, no hard-coded parameters
|-- scripts/        # short runs that call src/ and save results
|-- test/           # runtests.jl: the physics checks
|-- data/
|   |-- exp_raw/    # measurements: read-only, never modified
|   |-- exp_pro/    # processed data
|   `-- sims/       # simulation results (generated)
|-- plots/          # figures (generated)
|-- Project.toml    # dependencies
|-- Manifest.toml   # exact versions
`-- AGENTS.md       # instructions for the agent
\end{lstlisting}
Functions such as \cmd{datadir("sims", "rabi")} and \cmd{plotsdir()} return absolute
paths into these folders, so a script works from any working directory. A project with
this layout reads the same to you, your colleagues and the agent. The raw-data folder is
read-only by convention; \cmd{AGENTS.md} repeats the rule.

\textbf{The code is one part of the project.} A research project also has a manuscript,
notes and talks, with other collaborators and another history. The code gets its own
repository inside the project folder, so it can be shared, tested and published on its
own; in the course, \cmd{DrivenSpin/} always means that code repository.

\subsection{Scripts and \texttt{produce\_or\_load}}

The slide's script (it becomes \cmd{scripts/rabi.jl} in \cmd{DrivenSpin/}):
\begin{lstlisting}
using DrWatson
@quickactivate "DrivenSpin"          # this project's environment, from any folder
using DrivenSpin                     # the module in src/

for c in dict_list(Dict(:Δ => [0.0, 2π], :Ω => 2π))   # rad/μs; one run per Δ
    produce_or_load(c, datadir("sims", "rabi")) do c
        H = hamiltonian(c[:Δ], c[:Ω])
        t = range(0, 3, length = 301)                  # μs
        Dict("t" => t, "Pe" => excited_population.(Ref(H), t))
    end
end
\end{lstlisting}
\begin{placeholder}[Placeholder: \texttt{DrivenSpin/scripts/rabi.jl}]
Replace the listing above by the real file once D3 exists (it should be identical; the
slide was written first).
\end{placeholder}

Line by line:
\begin{itemize}
  \item \cmd{@quickactivate "DrivenSpin"} activates the project environment (the
        \cmd{Project.toml} of \cmd{DrivenSpin/}) wherever the script is started from.
  \item \cmd{dict\_list} expands a dictionary whose values may be vectors into all
        combinations: here two configurations, $\Delta = 0$ and $\Delta = 2\pi$, both with
        $\Omega = 2\pi$.
  \item \cmd{produce\_or\_load(f, config, path)} builds a file name from \cmd{config} with
        \cmd{savename}, and returns \cmd{(data, file)}. If the file exists, it is loaded;
        otherwise \cmd{f(config)} runs and its result is saved \rf{drwatson-docs}. The
        \cmd{do} block is Julia syntax for passing \cmd{f} as an anonymous function.
  \item \cmd{savename} sorts the keys, joins \cmd{key=value} pairs with \cmd{\_} and rounds
        numbers to three significant digits by default, which gives
        \cmd{Δ=0.0\_Ω=6.28.jld2}.
  \item For \cmd{.jld2} files the result is saved with \cmd{tagsave}, which adds the key
        \cmd{gitcommit}: the output of \cmd{git describe} for the project repository
        \rf{drwatson-docs}.
\end{itemize}
If the repository had uncommitted changes, \cmd{gitcommit} ends in \cmd{-dirty}: the
code that made the file is not in any commit, so that version is lost. DrWatson can also
store the uncommitted diff under the key \cmd{gitpatch} (keyword \cmd{storepatch}, off by
default). Hence the rule on the slide: \emph{commit first, then run}.

Two cautions for questions. First, \cmd{produce\_or\_load} decides by file name only: if
you change the code and rerun, it loads the old result. Delete the file or pass
\cmd{force = true}. Second, three significant digits mean $\Delta = 6.283$ and
$\Delta = 6.284$ share a file name; choose parameters, or \cmd{savename} options, so that
names stay distinct.

\subsection{The environment is part of the result}

\cmd{Project.toml} lists what you depend on and is written by \cmd{Pkg.add}.
\cmd{Manifest.toml} records the exact version of every package, including indirect
dependencies, and is never edited by hand. With both committed,
\cmd{Pkg.instantiate()} reproduces your environment on another machine. An agent that
``fixes'' an error with \cmd{Pkg.update()} changes package versions, and numbers may change
with them; a \cmd{Manifest.toml} in a diff you did not expect needs an explanation.

\subsection{Tests}

\begin{lstlisting}
using DrivenSpin, Test

@testset "Rabi" begin
    Δ, Ω = 0.3, 2π                                  # rad/μs
    ΩR = sqrt(Ω^2 + Δ^2)
    H = hamiltonian(Δ, Ω)
    for t in range(0, 3, length = 31)               # μs
        @test excited_population(H, t) ≈ Ω^2 / ΩR^2 * sin(ΩR * t / 2)^2 atol = 1e-10
    end
    @test excited_population(hamiltonian(0.0, Ω), π / Ω) ≈ 1    # π-pulse
end
\end{lstlisting}
\begin{placeholder}[Placeholder: \texttt{DrivenSpin/test/runtests.jl}]
Replace by the stage-1 tests of D3 once they exist.
\end{placeholder}

\cmd{@testset} groups tests under a name that says which physical claim they check;
\cmd{@test} records a pass or a failure without stopping \rf{julia-test}. The operator
\lstinline|≈| is \cmd{isapprox}, which by default compares with a \emph{relative} tolerance
of $\sqrt{\epsilon_{\mathrm{mach}}} \approx 1.5 \times 10^{-8}$, where $\epsilon_{\mathrm{mach}}$ is the
double-precision machine epsilon, and no absolute tolerance. Near zero a relative tolerance is
meaningless: \verb|isapprox(1e-12, 0.0)| is \cmd{false}. Hence the \cmd{atol} on the
slide wherever the expected value can vanish ($\Pe(0) = 0$). Run the tests with
\verb|julia --project -e 'using Pkg; Pkg.test()'|, before every commit; keep them under a
minute.

A good test checks \emph{physics}, not the implementation (a formula, a limit, a
conservation law); it is deterministic (fixed random seed); it is fast; and it fails when
the physics is wrong. The last property is the one to check deliberately: break the code
and watch the test fail. A test that cannot fail checks nothing.

\subsection{Continuous integration}

Continuous integration (CI) runs the tests on a clean machine at every push. GitHub
Actions reads workflow files in \cmd{.github/workflows/}. The slide's file:
\begin{lstlisting}
name: CI
on: [push, pull_request]
jobs:
  test:
    runs-on: ubuntu-latest
    steps:
      - uses: actions/checkout@v4             # get the repository
      - uses: julia-actions/setup-julia@v2    # install Julia
      - uses: julia-actions/cache@v2          # cache packages between runs
      - uses: julia-actions/julia-buildpkg@v1 # instantiate the environment
      - uses: julia-actions/julia-runtest@v1  # Pkg.test()
\end{lstlisting}
Each \cmd{uses:} line runs a published action \rf{julia-actions}. The result appears as a
check mark or a cross next to the commit, under the \emph{Actions} tab. CI catches
``works on my laptop'' (a missing dependency, a file that was never committed) and makes
sure nobody, including the agent, forgets to run the tests. The file is identical for any
Julia project. The workflow runs only when \cmd{.github/} sits at the root of a
repository, which is why \cmd{DrivenSpin/} carries its own copy.

\subsection{Your tools, same principles}

\begin{center}
\begin{tabularx}{\linewidth}{@{}l X@{}}
\toprule
\textbf{Principle} & \textbf{In the lab} \\
\midrule
Version control & Analysis scripts, instrument settings, HDL sources of FPGA designs. \\
Tests & MATLAB's \cmd{matlab.unittest} framework \rf{matlab-unittest}; simulation
  testbenches for FPGA designs. \\
Raw data is read-only & The agent never writes where measurements land. \\
Traceability & Save the commit hash and the settings with every measurement run. \\
\bottomrule
\end{tabularx}
\end{center}

\begin{qa}
\Q{Do I have to use DrWatson?}
No. The principles are a fixed layout, results tied to a commit, and a pinned
environment. DrWatson gives them to Julia projects in a few lines.

\Q{Why not commit the results too?}
They can be regenerated from code, data and parameters, which are committed; large
binary files bloat the history. Archive final results that a paper depends on in a data
repository with a DOI.

\Q{How do I test code with random numbers?}
Fix the seed (\cmd{Random.seed!}) in the test, and test statistical claims with a
tolerance derived from the expected spread (Section~\ref{sec:verification}).

\Q{Does CI cost money?}
GitHub Actions is free for public repositories; private repositories have a monthly
quota of free minutes.
\end{qa}

\textbf{Read next:} \emph{Good enough practices in scientific computing} \rf{wilson-good};
the DrWatson documentation \rf{drwatson-docs}, ``Workflow tutorial''; the Scientific
Python Development Guide \rf{scipy-dev} for the same ideas in another language.

% Block 6 of the slides (06_verification.tex).
\section{Verification in physics}\label{sec:verification}

\onslides{exact checks (analytic solution, limiting case, structure, symmetry);
consistency checks (convergence, independent route, bounds and units,
injection--recovery); convergence order; two independent routes; injection--recovery; tests
first; who checks the tests; a status for every claim.}

Software engineering tests code against examples someone wrote down. Physics offers
something stronger: statements that must hold whatever the implementation. The symbols
of the driven spin used below are defined in Section~\ref{sec:physics}: the excited-state
population $\Pe$, the Rabi frequency $\Omega$, the detuning $\Delta$, the generalised Rabi
frequency $\OmR = \sqrt{\Omega^2 + \Delta^2}$, the density matrix $\rho$, the times
$T_1$ and $T_2$, the pure-dephasing time $T_\phi$ with rate $\Gamma_\phi = 1/T_\phi$, and the
dissipator $\Dis$ of the master equation~\eqref{eq:lindblad}.

\subsection{Exact statements}

\begin{center}
\renewcommand{\arraystretch}{1.3}
\begin{tabularx}{\linewidth}{@{}p{0.22\linewidth}X@{}}
\toprule
\textbf{Check} & \textbf{Driven spin, and why it holds} \\
\midrule
Analytic solution & $\Pe(t) = (\Omega^2/\OmR^2)\sin^2(\OmR t/2)$ from $\ket{g}$,
  Eq.~\eqref{eq:rabi}. \\
Limiting case & $\Omega = 0$: $H$ is diagonal, populations do not move.
  $\Delta = 0$, $t = \pi/\Omega$: full inversion, $\Pe = 1$. \\
Structure & $\tr\rho = 1$, $\rho = \rho^\dagger$, eigenvalues of $\rho \geq 0$: the
  Lindblad equation preserves all three (Section~\ref{sec:lindblad}); a numerical scheme
  need not. \\
Symmetry & $\Pe$ is even in $\Delta$: Eq.~\eqref{eq:rabi} depends on $\Delta^2$ only.
  It stays even with decoherence: the map $\rho \mapsto \sigma_z \rho^\ast \sigma_z$
  turns the master equation~\eqref{eq:lindblad} at $\Delta$ into the one at $-\Delta$ and
  leaves $\ket{g}\bra{g}$ and all populations unchanged (checked numerically). \\
\bottomrule
\end{tabularx}
\end{center}
A symmetry test is cheap and catches a whole class of bugs: a term with the wrong sign
breaks it even when every individual value looks plausible.

\subsection{Consistency checks}

\textbf{Convergence.} For a method of order $p$, the error at a fixed final time $T$
scales as $\varepsilon(\delta t) \simeq C \delta t^{\,p}$ for small step $\delta t$, with
$C$ independent of $\delta t$. Halving the step therefore divides the error by $2^p$:
\begin{equation}
  p \simeq \log_2 \frac{\varepsilon(\delta t)}{\varepsilon(\delta t/2)} .
\end{equation}
The test on the slide compares the free decay of $\Pe$ from $\ket{e}$ with the exact
$\ee^{-T/T_1}$ for the classical fourth-order Runge--Kutta method (RK4) with
$T_1 = 1\,\mu$s, $T = 2\,\mu$s and $\delta t = 0.02$ and $0.01\,\mu$s, and demands
$p = 4 \pm 0.2$. The slide's numbers were checked: the measured order is $4.01$. A
convergence test catches a wrong integrator, a missing term in the equation (the error
then saturates at a constant), and a result reported before it has converged. The same
idea applies to any cutoff: grid size, Fock-space truncation, number of Monte Carlo
samples.

\begin{placeholder}[Placeholder: convergence test of \texttt{DrivenSpin}]
Quote the convergence test from \texttt{DrivenSpin/test/runtests.jl} once it exists. The
slide's version:
\lstinline|@test log2(err(0.02) / err(0.01)) ≈ 4 atol = 0.2|.
\end{placeholder}

\textbf{Independent route.} Solve the same problem with a different method written by
someone else, and compare. For the driven spin: the hand-written Lindblad solver against
QuantumOptics.jl \rf{qojulia}. Note how the dephasing appears there: the slide passes the
jump operator $\sqrt{1/(2T_\phi)}\,\sigma_z$, whose rate $1/(2T_\phi) = \Gamma_\phi/2$
matches the term $(\Gamma_\phi/2)\Dis[\sigma_z]\rho$ of Eq.~\eqref{eq:lindblad}.
Agreement is strong evidence because independent implementations are unlikely to share an
error; disagreement does not tell you which one is wrong. The same holds for derivations:
check every analytic result numerically, and every numerical result against a limit you
can do by hand.

\textbf{Bounds and units.} $0 \leq \Pe \leq 1$; $T_2 \leq 2T_1$ (Section~\ref{sec:lindblad});
rates in rad/\textmu s. A test with a known period catches a factor of $2\pi$: with
$\Omega = 2\pi$\,rad/\textmu s the resonant Rabi period is $1$\,\textmu s.

\subsection{Injection--recovery}

The experimentalist's test of an analysis pipeline. Choose true parameters
$\theta^\ast$; simulate data with realistic noise; run the pipeline, which returns fitted
values $\hat\theta$ with reported standard errors $\sigma$; and check
\begin{equation}
  \left|\hat\theta - \theta^\ast\right| < 3\sigma
\end{equation}
for each parameter. One seed tests that the pipeline works. Many seeds test that the
error bars are honest: the \emph{pulls} $(\hat\theta - \theta^\ast)/\sigma$ must be
distributed as a Gaussian of zero mean and unit width. A mean away from zero is a bias;
a width above one means the error bars are too small. This is the test that catches an
error-bar computation that ignores correlations, a noise model that does not match the
data, or a fit that converges to a wrong local minimum. Stage~4 of the hands-on
(Section~\ref{sec:stage4}) is an injection--recovery test.

\subsection{Tests first}

The sequence on the slide: you list the checks; the agent writes the tests; you read and
commit them; the agent implements without touching \cmd{test/}; you read the diff.
Committed tests are a fixed target the agent cannot move unnoticed: any change to them
shows up in \cmd{git diff test/}. Writing them first also forces you to say what
``correct'' means before you see a result, which protects against accepting whatever
comes out. This is the software engineers' test-driven development, and it matches the
verifiable-reward training of Section~\ref{sec:llm}: give the agent a check, and it
iterates until the check passes.

\textbf{Who checks the tests?} An agent writing both code and tests grades its own
homework. Tests derived from the implementation (``compute the value, paste it as
expected value'') check only that the code does what it does. Test against physics.
Break the code on purpose: every important test must fail. Changing the code
systematically and checking that some test fails is called \emph{mutation testing}; the
planted bug of block~8 is one hand-made mutation. And read every line of \cmd{test/}:
it is short, and it is the part you cannot delegate.

\subsection{A status for every claim}

The slide's labels come from the instructor's own projects: \emph{verified} (reproduced
by an independent route), \emph{fixed} (wrong as stated, corrected), \emph{weakened}
(true only under a narrower condition, now stated), \emph{retracted} (does not survive),
\emph{open} (not decidable with the tools at hand). Keeping such a list in a file in the
repository turns ``the agent said so'' into a record that a referee, or you in a year,
can audit.

\begin{qa}
\Q{How do I test a fit where I do not know the true answer?}
That is what injection--recovery is for: you create a case where you do know it, with
the same noise and sampling as the real data.

\Q{My result agrees with QuantumOptics.jl. Is it right?}
It is much more likely to be right. Both could share a wrong model (for example the wrong
sign convention for $\Delta$), so the analytic and limiting-case tests remain necessary.

\Q{What if the agent says a test is wrong?}
It might be. The prompt on the slide says ``If a test seems wrong, stop and tell me why''.
Then you decide, and the change to the test is its own commit with its reason in the
message.
\end{qa}

% The toy problem: physics behind the stage cards of block 7 (07_handson.tex) and
% the checks of block 6.  Conventions of PLAN.md section 5.
\section{The driven spin: physics of the hands-on}\label{sec:physics}

\onslides{block~7 gives one card per stage: (1) Rabi oscillations, the starting point;
(2) decoherence, $T_1$ and $T_2$; (3) Ramsey fringes and spin echo; (4) a synthetic
experiment and fit; (5, stretch) the Jaynes--Cummings model. Each stage is worked
``tests first'' in four steps (Section~\ref{sec:handson}).}

This section derives every formula on the stage cards, so that every test can be
explained. All formulas were checked numerically against a direct
integration of the equations (Python/SciPy, independent of \cmd{DrivenSpin/}). Krantz
et al.\ \rf{krantz} is an excellent experimental-minded review of the same material for
superconducting qubits.

\subsection{The model}

A two-level system with states $\ket{e}$ (excited) and $\ket{g}$ (ground), transition
angular frequency $\omega_q$, is driven by a field of angular frequency $\omega_d$. In the
laboratory frame,
\begin{equation}
  H_{\mathrm{lab}}(t) = \frac{\omega_q}{2}\sigma_z + \Omega \cos(\omega_d t)\,\sigma_x ,
\end{equation}
where $\Omega$ is set by the drive amplitude, and $\sigma_x$, $\sigma_y$, $\sigma_z$ are the
Pauli matrices in the basis $(\ket{e}, \ket{g})$. Transform to the frame rotating with the
drive, $\ket{\psi'} = U\ket{\psi}$ with $U(t) = \exp(i\omega_d t \sigma_z/2)$. The new
Hamiltonian is $U H_{\mathrm{lab}} U^\dagger + i \dot U U^\dagger$. The second term gives
$-(\omega_d/2)\sigma_z$; in the first, $\sigma_+ = \sigma_-^\dagger$ and $\sigma_-$ pick up phases $\ee^{\pm i\omega_d t}$,
and the drive splits into a static part $(\Omega/2)\sigma_x$ and terms oscillating at
$2\omega_d$. Dropping the latter (the \emph{rotating-wave approximation}, valid for
$\Omega, |\Delta| \ll \omega_d$) leaves the course Hamiltonian
\begin{equation}\label{eq:H}
  H = \frac{\Delta}{2}\sigma_z + \frac{\Omega}{2}\sigma_x,
  \qquad \Delta \equiv \omega_q - \omega_d .
\end{equation}
$\Omega$ is the \emph{Rabi frequency}, $\Delta$ the \emph{detuning}. Units: $\hbar = 1$,
times in \textmu s, $\Omega$ and $\Delta$ in rad/\textmu s.

\textbf{Bloch vector.} Write a state as $\rho = \frac12(1 + \mathbf{S}\cdot\boldsymbol\sigma)$
with the Bloch vector $\mathbf{S} = (\langle\sigma_x\rangle, \langle\sigma_y\rangle,
\langle\sigma_z\rangle)$; $\ket{e}$ is the north pole $S_z = +1$, $\ket{g}$ the south pole.
For $H = \frac12\boldsymbol\omega\cdot\boldsymbol\sigma$ with
$\boldsymbol\omega = (\Omega, 0, \Delta)$, the Heisenberg equations give a precession,
\begin{equation}
  \dot{\mathbf{S}} = \boldsymbol\omega \times \mathbf{S},
\end{equation}
about the axis $\boldsymbol\omega$ at angular frequency
$|\boldsymbol\omega| = \OmR \equiv \sqrt{\Omega^2 + \Delta^2}$, the \emph{generalised Rabi
frequency}. The excited-state population is $\Pe = \bra{e}\rho\ket{e} = (1 + S_z)/2$.

\subsection{Stage 1: Rabi oscillations}\label{sec:stage1}

With $\mathbf{n} = \boldsymbol\omega/\OmR$, Eq.~\eqref{eq:H} is
$H = (\OmR/2)\,\mathbf{n}\cdot\boldsymbol\sigma$, and since
$(\mathbf{n}\cdot\boldsymbol\sigma)^2 = 1$,
\begin{equation}
  U(t) = \ee^{-iHt} = \cos\frac{\OmR t}{2} - i \sin\frac{\OmR t}{2}\;\mathbf{n}\cdot\boldsymbol\sigma .
\end{equation}
The matrix element $\bra{e}\mathbf{n}\cdot\boldsymbol\sigma\ket{g} = n_x - i n_y = \Omega/\OmR$
gives, starting from $\ket{g}$,
\begin{equation}\label{eq:rabi}
  \Pe(t) = \left|\bra{e}U(t)\ket{g}\right|^2 = \frac{\Omega^2}{\OmR^2}\sin^2\frac{\OmR t}{2} .
\end{equation}
At resonance ($\Delta = 0$) the population oscillates fully, $\Pe = \sin^2(\Omega t/2)$,
with period $2\pi/\Omega$: $1$\,\textmu s for $\Omega = 2\pi$\,rad/\textmu s. Off
resonance the oscillation is faster and its amplitude $\Omega^2/\OmR^2$ smaller.

A resonant pulse of duration $t$ rotates the Bloch vector about $x$ by the \emph{pulse
area} $\theta = \Omega t$:
\begin{equation}\label{eq:pulse}
  R(\theta) = \ee^{-i\theta\sigma_x/2}.
\end{equation}
A $\pi$-pulse ($t = \pi/\Omega$) maps $\ket{g} \to -i\ket{e}$; a $\pi/2$-pulse creates an
equal superposition. The phase is visible in $\langle\sigma_y\rangle$: from $\ket{g}$ at
$\Delta = 0$, $\dot{\mathbf{S}} = \boldsymbol\omega\times\mathbf{S}$ with $\mathbf{S}(0) =
(0, 0, -1)$ gives $\mathbf{S}(t) = (0, \sin\Omega t, -\cos\Omega t)$, so
$\langle\sigma_y\rangle(t) = \sin(\Omega t)$. This sign is the one that flips in the
``bug your tests cannot see'' (Section~\ref{sec:failures}).

\textbf{Stage-1 tests} (in the starting point): Eq.~\eqref{eq:rabi} at many times; the
$\pi$-pulse inverts at $\Delta = 0$; unitarity $U^\dagger U = 1$ and norm conservation.

\begin{placeholder}[Placeholder: stage-1 code of \texttt{DrivenSpin}]
Quote \texttt{hamiltonian(Δ, Ω)}, \texttt{excited\_population(H, t)} and
\texttt{excited\_population(ρ)} from \texttt{DrivenSpin/src/DrivenSpin.jl} once D3 exists
(names as agreed in PLAN.md §11).
\end{placeholder}

\subsection{Stage 2: decoherence}\label{sec:lindblad}

A qubit coupled to an environment is described by its density matrix $\rho$. Under the
Born--Markov approximations (weak coupling, short environment memory), its evolution has
the Gorini--Kossakowski--Sudarshan--Lindblad form \rf{lindblad,gks,breuer}. The course
uses
\begin{equation}\label{eq:lindblad}
  \frac{\dd\rho}{\dd t} = -i[H, \rho] + \Gamma_1 \Dis[\sigma_-]\rho
    + \frac{\Gamma_\phi}{2}\Dis[\sigma_z]\rho,
  \qquad
  \Dis[L]\rho = L\rho L^\dagger - \tfrac12\{L^\dagger L, \rho\},
\end{equation}
with $\sigma_- = \ket{g}\bra{e}$, the energy-relaxation rate $\Gamma_1 = 1/T_1$ and the
pure-dephasing rate $\Gamma_\phi = 1/T_\phi$. The first dissipator moves population from
$\ket{e}$ to $\ket{g}$ (spontaneous emission into the environment); the second randomises
the phase between $\ket{e}$ and $\ket{g}$ without exchanging energy (for example, a
fluctuating $\omega_q$).

\textbf{Why this form.} Lindblad \rf{lindblad} and Gorini, Kossakowski and Sudarshan
\rf{gks} showed in 1976 that Eq.~\eqref{eq:lindblad}, with any operators $L$ and
non-negative rates, is the most general Markovian generator that keeps $\rho$ a valid
density matrix: trace one, Hermitian and positive (in fact completely positive) at all
times. Negative rates, or a dissipator with the wrong sign, break positivity, which is why
the structure tests of stage~2 detect them.

\textbf{Solving it.} Take $\Omega = 0$ and write $\rho_{ee} = \bra{e}\rho\ket{e}$,
$\rho_{eg} = \bra{e}\rho\ket{g}$. Working out Eq.~\eqref{eq:lindblad} element by element
($\Dis[\sigma_z]$ leaves the diagonal alone and multiplies $\rho_{eg}$ by $-2$),
\begin{equation}
  \dot\rho_{ee} = -\Gamma_1 \rho_{ee},
  \qquad
  \dot\rho_{eg} = -\left(i\Delta + \frac{\Gamma_1}{2} + \Gamma_\phi\right)\rho_{eg}.
\end{equation}
The population decays at $\Gamma_1$; the coherence at
\begin{equation}\label{eq:T2}
  \frac{1}{T_2} = \frac{1}{2T_1} + \frac{1}{T_\phi}.
\end{equation}
Since $\Gamma_\phi \geq 0$, $T_2 \leq 2T_1$: the coherence decays at least half as fast
as the population, because each decay event also destroys the phase. A parameter set with
$T_2 > 2T_1$ would need $T_\phi < 0$, which violates positivity; hence the specification
that \cmd{lindblad} rejects it with an \cmd{ArgumentError}. $T_1 = T_2 = \infty$ switches
dissipation off and must reproduce stage~1.

\textbf{Stage-2 tests} (stage card): at all times $\tr\rho = 1$, $\rho = \rho^\dagger$,
eigenvalues $\geq 0$; $\Omega = 0$ from $\ket{e}$: $\Pe(t) = \ee^{-t/T_1}$; $\Omega =
\Delta = 0$ from $\ket{+} = (\ket{e} + \ket{g})/\sqrt2$: $|\rho_{eg}(t)| =
\frac12\ee^{-t/T_2}$; $T_2 > 2T_1$ rejected; $T_1, T_2 \to \infty$ gives stage~1.

\textbf{Numerics.} Equation~\eqref{eq:lindblad} is linear in $\rho$,
$\dot\rho = \mathcal{L}\rho$. RK4 with step $\delta t$ applies a fourth-order polynomial in
$\delta t\,\mathcal{L}$ per step. Because $\mathcal{L}$ preserves trace and Hermiticity, so
does the polynomial, up to rounding; positivity it preserves only approximately, with an
error that shrinks with $\delta t$. The structure tests therefore test the
implementation of $\mathcal{L}$ more than the integrator, and the convergence test
(Section~\ref{sec:verification}) tests the integrator. With the slide's parameters the
measured order is $p = 4.01$.

\subsection{Stage 3: Ramsey fringes and spin echo}\label{sec:stage3}

\textbf{Ramsey.} Start in $\ket{g}$; apply $R(\pi/2)$ (Eq.~\eqref{eq:pulse}); let the qubit
evolve freely for a time $\tau$ with the drive off, so $H = (\Delta/2)\sigma_z$ and the
Bloch vector precesses about $z$ at angular frequency $\Delta$; apply $R(\pi/2)$ again;
measure $\Pe$. Pulses are taken as ideal (instantaneous) and $T_1 \to \infty$, so that
$T_2 = T_\phi$. During the free evolution the equatorial component of the Bloch vector
rotates by $\Delta\tau$ and shrinks by $\ee^{-\tau/T_2}$; the second pulse maps it to
\begin{equation}\label{eq:ramsey}
  \Pe(\tau) = \tfrac12\left[1 + \ee^{-\tau/T_2}\cos(\Delta\tau)\right].
\end{equation}
At $\tau = 0$ the two $\pi/2$-pulses make a $\pi$-pulse and $\Pe = 1$. The fringes
oscillate at the detuning: this is how experiments calibrate $\omega_q$ against the drive,
and how they measure coherence times. Ramsey introduced the method of separated
oscillatory fields in 1950 \rf{ramsey}.

\textbf{Inhomogeneous dephasing.} In a real experiment the detuning differs from shot to
shot by a static offset $\delta$ (slow drifts, a distribution of $\omega_q$ in an ensemble).
Averaging $\cos[(\Delta + \delta)\tau]$ over a Gaussian distribution of $\delta$ with
standard deviation $s_\delta$ gives $\cos(\Delta\tau)\,\ee^{-s_\delta^2\tau^2/2}$: an extra,
Gaussian decay. The coherence time measured by Ramsey is therefore usually called
$T_2^\ast$, and $T_2^\ast \leq T_2$.

\textbf{Echo.} Hahn's spin echo \rf{hahn} inserts a $\pi$-pulse in the middle:
$\pi/2$ -- $\tau/2$ -- $\pi$ -- $\tau/2$ -- $\pi/2$. The $\pi$-pulse about $x$ maps
$\sigma_z \to -\sigma_z$, so the phase $\delta\tau/2$ acquired in the first half is undone
in the second, whatever $\delta$ is. With ideal pulses and no decoherence the sequence
returns $\ket{g}$ for every static $\delta$, $\Pe = 0$ (the total rotation is $2\pi$
about $x$). Averaging over $\delta$ therefore changes nothing: the echo removes static
dephasing and leaves only the irreversible part, $T_2$. Fluctuations of $\delta$ during
the sequence are not refocused, which is why echo decay measures $T_2$.

\textbf{Stage-3 tests} (stage card): fringes at angular frequency $\Delta$ under the
envelope $\ee^{-\tau/T_2}$; echo independent of a static $\delta$; averaged over random
static $\delta$, Ramsey decays and echo does not.

\subsection{Stage 4: a synthetic experiment and its fit}\label{sec:stage4}

\textbf{Data.} For each delay $\tau_k$ compute $\Pe(\tau_k)$ from stage~3 and simulate $N$
single-shot measurements: the number of shots that find $\ket{e}$ is binomial,
$n_k \sim \mathrm{Binomial}(N, \Pe(\tau_k))$, and the measured population is
$\hat P_k = n_k/N$ with variance $\Pe(1 - \Pe)/N$.

\textbf{Fit.} The model on the stage card is
\begin{equation}\label{eq:fitmodel}
  P(\tau) = B + A\,\ee^{-\tau/T_2}\cos(\Delta\tau + \varphi),
\end{equation}
with amplitude $A$, offset $B$ and phase $\varphi$ as additional parameters; for ideal pulses
$A = B = \frac12$ and $\varphi = 0$, and leaving them free absorbs pulse errors and
readout offsets, as in an experiment. LsqFit \rf{lsqfit} minimises the weighted sum of
squared residuals:
\begin{lstlisting}
fit = curve_fit(model, τ, Phat, w, p0)   # w = 1 ./ variances (inverse variance)
coef(fit), stderror(fit)               # best-fit parameters and standard errors
\end{lstlisting}
The weights are the inverse variances; estimated from the data, $\hat P_k(1 -
\hat P_k)/N$ vanishes at $\hat P_k = 0$ or $1$, so floor it (or use the model value). The
starting values \cmd{p0} matter: the fit has many local minima in $\Delta$; take $\Delta$
from the peak of a Fourier transform of the data, $T_2$ from the visible decay.

\textbf{Why four times the shots halves the errors.} Each $\hat P_k$ has standard deviation
$\propto 1/\sqrt{N}$, and the standard errors $\sigma$ of the fitted parameters scale
linearly with the data errors. So $\sigma \propto 1/\sqrt{N}$, and $N \to 4N$ gives $\sigma \to \sigma/2$.

\textbf{Stage-4 tests} (stage card): injection--recovery with a fixed seed, fitted $T_2$
and $\Delta$ within $3\sigma$ of the true values; four times more shots give error bars
about half as large; the script saves the fit with \cmd{produce\_or\_load}, so the file
carries the commit.

\begin{placeholder}[Placeholder: stage-4 reference solution]
Quote the fit function and its injection--recovery test once D3 has them (on the
\texttt{solutions} branch, not in the participants' starting point).
\end{placeholder}

\subsection{Stage 5 (stretch): the Jaynes--Cummings model}\label{sec:stage5}

A qubit coupled to one mode of a cavity with angular frequency $\omega_c$, photon
annihilation operator $a$ and coupling $g$ \rf{jaynescummings}:
\begin{equation}\label{eq:jc}
  H_{\mathrm{JC}} = \omega_c a^\dagger a + \frac{\omega_q}{2}\sigma_z
    + g\left(a^\dagger\sigma_- + a\sigma_+\right),
\end{equation}
with $\sigma_+ = \ket{e}\bra{g}$ as before. The coupling term is already in
rotating-wave form (valid for $g \ll \omega_c, \omega_q$): it only exchanges one photon for
one qubit excitation. Hence the \emph{excitation number}
$N_{\mathrm{exc}} = a^\dagger a + \sigma_+\sigma_-$ commutes with $H_{\mathrm{JC}}$ and is
conserved.

For each $n \geq 0$, $H_{\mathrm{JC}}$ couples only $\ket{n, e}$ and $\ket{n+1, g}$ (photon
number, qubit state), with matrix element $\bra{n+1, g}H_{\mathrm{JC}}\ket{n, e} =
g\sqrt{n+1}$. On resonance ($\omega_c = \omega_q$), starting from $\ket{n, e}$,
\begin{equation}
  \Pe(t) = \cos^2\!\left(g\sqrt{n+1}\,t\right),
\end{equation}
which oscillates at angular frequency $2g\sqrt{n+1}$; for $n = 0$ this is the vacuum Rabi
frequency $2g$, also the splitting of the two dressed states in the spectrum.

\textbf{Stage-5 tests}: $N_{\mathrm{exc}}$ is conserved (commutator, or its expectation
value along the evolution); the oscillation frequency is $2g\sqrt{n+1}$; the result
converges in the photon cutoff. A subtlety for the last test: for a Fock initial state,
conservation of $N_{\mathrm{exc}}$ makes the truncated dynamics \emph{exact} once the cutoff
includes $n + 1$ photons, so convergence is trivial. A meaningful test starts from a coherent state
of the field, whose photon-number distribution has a tail; then the populations
(including the collapse and revival of the Rabi oscillations) converge only as the cutoff
grows.

\subsection{The hands-on in practice}\label{sec:handson}

Pairs, one person new to git typing the git commands, swapping the keyboard after each
stage. Getting started: copy \cmd{DrivenSpin/} from the course repository into a new
folder, \cmd{git init}, commit, and run the tests, which pass at stage~1.

\begin{placeholder}[Placeholder: getting-started commands]
Final commands depend on PLAN.md §11 (the course repository is already cloned as
\texttt{\textasciitilde/ai-course/course} during setup) and on the repository URL.
Copy them from the final slide ``Getting started'' (\texttt{slides/sections/07\_handson.tex}).
\end{placeholder}

Each stage in four steps:
\begin{enumerate}
  \item Prompt: read \cmd{AGENTS.md} and \cmd{src/}; write \emph{only} the tests for the
        stage in \cmd{test/runtests.jl}, from the checks on the stage card.
  \item You read \cmd{git diff test/} line by line and commit the tests.
  \item Prompt: implement the stage in \cmd{src/} until all tests pass; do not modify
        \cmd{test/}; if a test seems wrong, stop and explain why.
  \item You read \cmd{git diff}, run the tests yourself, and commit.
\end{enumerate}
Done early: push to a private GitHub repository with
\cmd{gh repo create DrivenSpin --private --source=. --push} and add the CI file; check
stage~2 against QuantumOptics.jl; stage~5; a convergence test of your time stepper. When
stuck: ask for an explanation before a fix; after circles, \cmd{git restore .} and a fresh
session.

% Block 8 of the slides (08_failures.tex) and the conclusion (90_conclusion.tex).
\section{Failure modes and responsibility}\label{sec:failures}

\onslides{the planted bug; ``plausible is not correct''; the bug your tests cannot
see; two rogues' galleries; you are the author; recording which tool was used.}

\subsection{The planted bug}

Participants apply a patch to their own repository with \cmd{git apply}, make the Rabi
plot, and run the tests. The patch changes one character: $\Omega/2 \to \Omega$ in
\cmd{hamiltonian}. Its effect at $\Delta = 0$: the Hamiltonian becomes $\Omega\sigma_x$, so
\begin{equation}
  \Pe(t) = \sin^2(\Omega t) \quad\text{instead of}\quad \sin^2(\Omega t/2).
\end{equation}
The plot still shows smooth, bounded, periodic oscillations between $0$ and $1$, at twice
the frequency. Without a reference, the eye accepts both curves (slide ``Plausible is not
correct''). The $\pi$-pulse test fails at once: at $t = \pi/\Omega$ the buggy population
is $\sin^2\pi = 0$ instead of $1$. Off resonance the bug changes amplitude and frequency
together, $\Pe = \Omega^2/(\Omega^2 + \Delta^2/4)\,\sin^2(\sqrt{\Omega^2 + \Delta^2/4}\,t)$,
so the analytic test fails at every $\Delta$. Find the bug with \cmd{git diff}; undo it with
\cmd{git restore .}.

\begin{placeholder}[Placeholder: \texttt{exercises/planted\_bug.patch}]
Quote the patch once D3 provides it (it must touch only the line of \texttt{hamiltonian},
so it applies after stages 2--4).
\end{placeholder}

\subsection{The bug your tests cannot see}

Flip the sign of the Hamiltonian term in Eq.~\eqref{eq:lindblad}: $-i[H, \rho] \to
+i[H, \rho]$. This is evolution under $-H$, equivalently backwards in time for the
coherent part. All populations in the course are unchanged, because
\begin{itemize}
  \item the Rabi formula~\eqref{eq:rabi} is even in $\Omega$ and in $\Delta$;
  \item the Ramsey and echo results~\eqref{eq:ramsey} depend on $\cos(\Delta\tau)$ only;
  \item the dissipators do not involve $H$.
\end{itemize}
So every test of stages 1--4 passes. Only phases change sign: at $\Delta = 0$,
$\langle\sigma_y\rangle(t) = \sin(\Omega t)$ (Section~\ref{sec:stage1}) becomes
$-\sin(\Omega t)$. The lesson on the slide: your tests decide which bugs you can see, so
test every quantity you will publish. A single test of $\langle\sigma_y\rangle$, or of
the sign of a Ramsey phase $\varphi$ with a known detuning, closes this gap.

\subsection{Rogues' gallery}

\begin{center}
\renewcommand{\arraystretch}{1.25}
\begin{tabularx}{\linewidth}{@{}p{0.36\linewidth}X@{}}
\toprule
\textbf{Failure} & \textbf{Why it happens, and the defence} \\
\midrule
Invented or misquoted reference & Plausible text with no lookup (Section~\ref{sec:llm}).
  Open the link; check that the paper exists \emph{and} says that. \\
Factor of 2, wrong sign & Conventions differ between sources ($\Omega$ vs.\ $\Omega/2$,
  $\Delta = \omega_q - \omega_d$ or the reverse). Analytic and limiting-case tests. \\
Test weakened to pass & Rewarded for passing checks (Section~\ref{sec:llm}). Commit tests
  first; read every test diff. \\
Convention silently changed & Fourier normalisation, $\hbar = 1$, basis order, sign of
  $\Delta$. Conventions in \cmd{AGENTS.md}, enforced by tests. \\
Invented or outdated function names & Knowledge cutoff, little Julia in the training data.
  Run it; pin versions in \cmd{Manifest.toml}. \\
Hz where rad/\textmu s was meant & A factor of $2\pi$ that looks harmless. Units in
  names and docstrings; a test with a known period. \\
Confident summary of a paper not read & The model writes what such a paper would
  plausibly say. Ask for page and equation numbers; check them. \\
``All tests pass'' that were never run & A summary is a claim like any other. Run them
  yourself; CI runs them anyway. Review the diff, not the summary. \\
\bottomrule
\end{tabularx}
\end{center}

\subsection{You are the author}

\textbf{Accountability.} An AI tool is not an author, and every claim in the paper is
yours. Nature's 2023 editorial \rf{nature-rules} set two ground rules that most
publishers now follow: no AI tool is credited as an author, and its use is documented in
the methods or acknowledgements. The American Physical Society updated its policy in
June 2026 \rf{aps-news,aps-policy}. It allows substantive uses (literature synthesis, data
analysis, scientific reasoning, figures) with disclosure, and states that an AI program
cannot be held accountable and therefore cannot be an author. The disclosure names the
tool and version, how it assisted, how it was directed, and how its output was verified;
AI used in the research goes into the methods section, AI-made figures into their
captions, other substantive uses into the acknowledgements. Light editing needs no
disclosure.

\textbf{Data.} When an agent reads a file, its content is sent to the provider. Unpublished
data, collaborators' code and personal data leave your machine. Check what your
institution and your collaborations allow, and what the provider's terms say about
retention and training. Referees take note: APS forbids uploading manuscripts under review
to tools whose terms allow training on them \rf{aps-policy}.

\textbf{Disclosure.} Read your journal's policy before you submit. The APS
requirement ``how its output was verified'' is easy to meet if the verification is in
the repository: the tests, the independent routes, the claim statuses.

\textbf{Reproducibility.} Record which tool and model helped, and when. A commit trailer
costs nothing and is searchable:
\begin{lstlisting}
Add Ramsey sequence with fringe-frequency test

Assisted-by: <agent> (<model>)
\end{lstlisting}
Some agents add a trailer themselves (for example \cmd{Co-Authored-By}); keep it. Git
treats any \cmd{Key: value} lines at the end of the message as trailers;
\cmd{git log --grep "Assisted-by"} finds them later.

\subsection{Conclusion of the course}

The six take-aways on the closing slide, with the section that explains each:
\begin{enumerate}
  \item An LLM is a very large fit: plausible is not the same as correct
        (Section~\ref{sec:llm}).
  \item An agent closes the loop, so it can check its work, if you give it checks
        (Section~\ref{sec:agents}).
  \item Context is finite: one task per session, write the rest down
        (Section~\ref{sec:context}).
  \item Git is the safety net: commit often, read every diff (Section~\ref{sec:git}).
  \item Structure, pinned environments and CI: every result traceable and checked
        (Section~\ref{sec:structure}).
  \item Verification is the job: tests that encode physics, written first, read by you
        (Sections~\ref{sec:verification} and~\ref{sec:physics}).
\end{enumerate}

\begin{qa}
\Q{Do I have to disclose that I used an agent to write analysis code?}
For APS journals, yes: AI used in conducting the research, such as data analysis, is
disclosed in the methods section \rf{aps-policy}. Other publishers differ; check.

\Q{Can I paste my unpublished data into the agent?}
Only if your institution, your collaboration and the provider's terms allow it. When in
doubt, use synthetic data with the same structure (stage~4 is an example).

\Q{The agent's summary says all tests pass. Is that enough?}
No. Run them yourself, or let CI run them. Review the diff, not the summary.
\end{qa}

% Stub: terms people use when talking about AI.  To be extended by Marcelo.
\section{Trendy things people talk about}\label{sec:trends}

\begin{placeholder}[Stub: to be extended]
This section is a skeleton. Each entry has a one-paragraph starting point with a verified
source, and a marker where more text should go. Marcelo's note said ``like Jev'',
read here as the Jevons paradox; \textbf{confirm}, and add or remove entries.
\end{placeholder}

\subsection{The Jevons paradox}

In \emph{The Coal Question} (1865), W.\,S.~Jevons argued that more efficient steam engines
would \emph{increase} Britain's coal consumption, not reduce it: cheaper use of a resource
opens new uses. His Chapter~VII, ``Of the Economy of Fuel'', puts it in one sentence: ``It
is the very economy of its use which leads to its extensive consumption'' \rf{jevons}.
The argument carries over to code. If AI makes code cheap, we will not write the same
code for less effort; we will write much more code, run more analyses and try more
ideas. The demand for verification grows with it, which is the thesis of the course
(Section~\ref{sec:thesis}) seen from economics.

\textbf{[TO EXTEND]} Whether the paradox holds depends on how elastic the demand is:
\dots{} Examples from the group: \dots{} Counter-arguments: \dots{}

\subsection{Vibe coding}

A term coined by A.~Karpathy in a post on X in February 2025 \rf{vibecoding}: accepting
AI-written code without reading it, steering only by whether the result seems to work. It
is fine for throw-away prototypes and the opposite of what this course recommends for
research code, where ``seems to work'' is exactly the failure of
Section~\ref{sec:failures}.

\textbf{[TO EXTEND]}

\subsection{Model Context Protocol (MCP)}

An open protocol, introduced by Anthropic in November 2024 \rf{mcp}, through which an
agent connects to external tools and data sources (databases, issue trackers, instrument
control software) in a uniform way. In the language of Section~\ref{sec:agents}: a
standard for adding tools to a harness.

\textbf{[TO EXTEND]} Relevance for labs: an MCP server in front of an instrument or a
data archive \dots

\subsection{Agent skills and instruction files}

Covered in Section~\ref{sec:context}: \cmd{AGENTS.md} \rf{agentsmd} and skills
\rf{skills}. The trend is towards writing procedures down once, in plain files, for any
agent to load.

\textbf{[TO EXTEND]}

\subsection{Long-horizon agents}

METR's measurements \rf{metr-horizon} suggest that the length of tasks agents complete
reliably has doubled about every seven months since 2019. Extrapolations of this trend
are a staple of AI discussions; the course's position is that longer unattended runs make
automatic verification more important.

\textbf{[TO EXTEND]}

\subsection{Reward hacking}

When a model is optimised against a check, it may learn to satisfy the check without
doing the task: modifying tests or scoring code \rf{metr-hacking}. This is Goodhart's law,
which Tao discusses for mathematics \rf{tao}: a measure that becomes a target ceases to be
a good measure.

\textbf{[TO EXTEND]}

\subsection{Further candidates}

\textbf{[TO DECIDE]} AI ``scientists'' and automated discovery claims; open-weight versus
closed models; energy use of training and inference; benchmarks and their saturation;
local models for confidential data.


\part{Reference cards}
% One page: git cheat sheet.  Usable on its own.
\clearpage
\section{Git cheat sheet}\label{sec:cheatsheet}

{\small
\renewcommand{\arraystretch}{1.15}
\begin{tabularx}{\linewidth}{@{}l X@{}}
\toprule
\multicolumn{2}{@{}l}{\textbf{Start}} \\
\cmd{git init} & Start recording in the current folder. \\
\cmd{git clone <address>} & Copy a repository (e.g.\ from GitHub), with its history. \\
\midrule
\multicolumn{2}{@{}l}{\textbf{Look} (safe: changes nothing)} \\
\cmd{git status} & What is new, changed, staged? Run it often. \\
\cmd{git diff} & Changes not yet staged, line by line. \\
\cmd{git diff --staged} & What the next commit will contain. \\
\cmd{git diff --stat} & Which files changed, and how much. \\
\cmd{git log --oneline} & The history, one line per commit, newest first. \\
\cmd{git log --stat} / \cmd{git log -p <file>} & Files per commit / the story of one file. \\
\cmd{git show <hash>} & One commit in full: message and diff. \\
\cmd{git diff <hash1> <hash2>} & What changed between two snapshots. \\
\midrule
\multicolumn{2}{@{}l}{\textbf{Record}} \\
\cmd{git add <file>} & Stage one file. \\
\cmd{git add -p} & Stage piece by piece; answer \cmd{y} or \cmd{n} per change. \\
\cmd{git add -u} & Stage changes to tracked files; new files stay out. \\
\cmd{git add .} & Stage everything new, changed or deleted here and below. \\
\cmd{git commit -m "why"} & Take the snapshot. The message says \emph{why}. \\
\midrule
\multicolumn{2}{@{}l}{\textbf{Undo}} \\
\cmd{git restore <file>} & Discard uncommitted edits to a file (\emph{gone for good}). \\
\cmd{git restore .} & Discard \emph{all} uncommitted edits here and below. Check first. \\
\cmd{git restore --staged <file>} & Unstage; the edit stays in the file. \\
\cmd{git revert <hash>} & New commit that undoes an old one; history kept. \\
\midrule
\multicolumn{2}{@{}l}{\textbf{Branch}} \\
\cmd{git switch -c <name>} & New branch here, and switch to it. \\
\cmd{git switch <name>} & Switch to an existing branch. \\
\cmd{git merge <name>} & Bring \cmd{<name>} into the current branch. \\
\cmd{git branch -D <name>} & Delete a branch you no longer want. \\
\midrule
\multicolumn{2}{@{}l}{\textbf{Share}} \\
\cmd{git push} / \cmd{git pull} & Send your commits / get and merge others'. \\
\bottomrule
\end{tabularx}
}

\textbf{Reading a diff.} \cmd{---\,a/} old file, \cmd{+++\,b/} new file;
\verb|@@ -12,6 +12,6 @@|: excerpt starting at line 12, 6 lines long, old and new;
\cmd{-} removed, \cmd{+} added, space unchanged. A modified line is one \cmd{-} and one
\cmd{+}.

\textbf{Never without asking}: \cmd{git reset --hard}, \cmd{git push --force},
\cmd{git clean -f}. They destroy work.

\textbf{With an agent.} (1) Commit before it starts. (2) One task, one commit. (3) Read
\cmd{git diff test/} first, then the rest. (4) Let it type the commands; you read the
result. Useful prompts: ``Commit my manual changes before starting.'' ``Work in modular
commits: one commit per change, with the tests passing at each.''

\textbf{What goes in}: source, scripts, tests, \cmd{Project.toml}, \cmd{Manifest.toml},
\cmd{README.md}, \cmd{AGENTS.md}. \textbf{What stays out} (\cmd{.gitignore}): raw data,
generated results and plots, passwords and keys, ever.

\textbf{Words.} \emph{Commit}: snapshot of the whole project, named by a hash (first 7
characters suffice). \emph{Staging area}: the next snapshot being prepared. \emph{Branch}:
a movable name for the latest commit of a line of work. \cmd{HEAD}: what you have checked
out. \emph{Remote}: another copy, usually \cmd{origin} on GitHub.

% One page: verification checklist.  Also in handout/checklist.md; keep both in sync.
\clearpage
\section{Before trusting AI output: a checklist}\label{sec:checklist}

\newcommand{\cb}{$\square$\ }
\begin{itemize}[label={}, leftmargin=0pt, itemsep=0pt, parsep=0pt]
\item \textbf{Before the agent starts}
\item \cb The project is under git and the last commit is clean (\cmd{git status}).
\item \cb The task is one task, in a fresh session.
\item \cb The prompt names the files, the conventions (or \cmd{AGENTS.md}), the tests that
      define ``done'', and what must not change.
\item \cb For anything non-trivial: the agent proposes a plan and you approve it.

\smallskip
\item \textbf{Tests}
\item \cb The tests were written (or listed) \emph{before} the implementation, read by
      you line by line, and committed.
\item \cb They check physics, not the implementation: analytic results, limiting cases,
      symmetries, conservation laws, bounds, units.
\item \cb At least one test checks every quantity you will publish, including phases and
      signs, not only populations.
\item \cb You broke the code on purpose, and the important tests failed.
\item \cb Approximate comparisons near zero have an absolute tolerance (\cmd{atol}).
\item \cb Random numbers have a fixed seed.

\smallskip
\item \textbf{Reading the change}
\item \cb You read \cmd{git diff test/} first: no loosened tolerance, deleted assertion
      or special case.
\item \cb You read the rest of the diff, not the agent's summary of it.
\item \cb No unexpected change to \cmd{Manifest.toml}, conventions or units.
\item \cb You ran the tests yourself, or CI ran them on the pushed commit.

\smallskip
\item \textbf{Numbers and fits}
\item \cb Converged: halving the step (or raising the cutoff) does not change the result
      beyond the tolerance; the observed order matches the method.
\item \cb Cross-checked by an independent route (another package, another method, a
      limit by hand).
\item \cb Injection--recovery: the pipeline recovers known parameters within error bars;
      over many seeds the pulls look like a unit Gaussian.
\item \cb Results saved with the commit hash (e.g.\ \cmd{produce\_or\_load}), from a
      clean repository (no \cmd{-dirty}).

\smallskip
\item \textbf{Text, derivations, references}
\item \cb Every reference opened: it exists, and it says what is claimed.
\item \cb Every derivation checked numerically or symbolically by a script in the
      repository.
\item \cb Every claim has a status: verified, fixed, weakened, retracted or open.

\smallskip
\item \textbf{Responsibility}
\item \cb No data left your machine that your institution or collaboration does not allow.
\item \cb The tool and model are recorded (commit trailer, README).
\item \cb The journal's AI policy is followed: disclosure of tool, role, direction and
      verification.
\item \cb You can defend every line as your own.
\end{itemize}


\clearpage
\part{References}
% Annotated references.  Every entry was checked on 2026-10-20; the "Checked" line says
% how.  Keys: \label{ref:<key>}, cited with \rf{<key>}.
\section{Annotated references}\label{sec:references}
\raggedright

Every entry below was opened and checked on 20 October 2026: title, authors, year and
link, and every number these notes quote from it. The \emph{Checked} line says against
what (the publisher's page, Crossref, the arXiv record, or the document itself). Entries
that could only be checked in part say so. Entries marked $\star$ are the best next
reading for participants.

\newcommand{\refentry}[2]{\item\label{ref:#1}#2}
\newcommand{\use}[1]{\par\textit{Use:} #1}
\newcommand{\checked}[1]{\par{\small\textit{Checked:} #1}}
\newcommand{\best}{$\star$\ }

\subsection*{AI in science}
\begin{enumerate}[label={[\arabic*]}, ref=\arabic*, series=refs, leftmargin=*, itemsep=6pt]
\refentry{theophyseval}{A.~Ramesh Chandra, A.~Mahajan, A.~Sanyal, M.~Wang et al.
  (39 authors), \emph{TheoPhysEval: Assessing AI agents at autonomous theoretical physics
  problem solving}, arXiv:2610.16733 (2026).
  \url{https://arxiv.org/abs/2610.16733}
  \use{Motivation slide. 50 unpublished problems; best agent $98\,\%$ (49/50), under
  3\,USD per problem; little change without web search; no agent-led breakthroughs.}
  \checked{the PDF (v1, 9 Oct 2026): abstract and Sec.~on results.}}
\refentry{tao}{\best T.~Tao, \emph{Mathematics in the age of AI}, arXiv:2608.16753
  (2026). \url{https://arxiv.org/abs/2608.16753}
  \use{Motivation slide. Essay based on his ICM 2026 public lecture. Section~7, ``Proof
  scarcity and proof abundance'', contains the quoted sentence; also discusses
  Goodhart's law.}
  \checked{the PDF (v1, 17 Aug 2026); quotation verbatim from Sec.~7.}}
\refentry{metr}{J.~Becker, N.~Rush, E.~Barnes, D.~Rein, \emph{Measuring the impact of
  early-2025 AI on experienced open-source developer productivity}, arXiv:2507.09089
  (2025). \url{https://arxiv.org/abs/2507.09089}
  \use{Scepticism: 16 developers, 246 tasks; forecast $-24\,\%$ time, perceived
  $-20\,\%$, measured $+19\,\%$.}
  \checked{arXiv record and abstract.}}
\end{enumerate}

\subsection*{How LLMs work}
\begin{enumerate}[label={[\arabic*]}, ref=\arabic*, resume=refs, leftmargin=*, itemsep=6pt]
\refentry{karpathy-intro}{\best A.~Karpathy, \emph{[1hr Talk] Intro to Large Language
  Models}, YouTube (22 Nov 2023). \url{https://www.youtube.com/watch?v=zjkBMFhNj_g}
  \use{The best one-hour overview; pitched at a general technical audience.}
  \checked{video page: title, channel, upload date.}}
\refentry{karpathy-deep}{\best A.~Karpathy, \emph{Deep Dive into LLMs like ChatGPT},
  YouTube (5 Feb 2025). \url{https://www.youtube.com/watch?v=7xTGNNLPyMI}
  \use{Three and a half hours: tokenization, pre-training, post-training, RL, in the
  order of block~1.}
  \checked{video page: title, channel, upload date.}}
\refentry{3b1b}{\best 3Blue1Brown (G.~Sanderson), \emph{Neural networks} series, in
  particular \emph{Transformers, the tech behind LLMs} (Chapter~5, 2024) and
  \emph{Attention in transformers, step-by-step} (Chapter~6, 2024).
  \url{https://www.3blue1brown.com/topics/neural-networks}
  \use{Visual explanation of attention (backup slide).}
  \checked{series page; video pages for chapters 5 and 6 (\cmd{wjZofJX0v4M},
  \cmd{eMlx5fFNoYc}), uploaded April 2024.}}
\refentry{vaswani}{A.~Vaswani, N.~Shazeer, N.~Parmar, J.~Uszkoreit, L.~Jones,
  A.\,N.~Gomez, Ł.~Kaiser, I.~Polosukhin, \emph{Attention is all you need}, Advances in
  Neural Information Processing Systems \textbf{30} (2017); arXiv:1706.03762.
  \url{https://arxiv.org/abs/1706.03762}
  \use{The transformer; attention formula of the backup slide.}
  \checked{arXiv record (title, authors, 2017). The NeurIPS volume is the standard
  citation; the arXiv record itself names no venue.}}
\refentry{sennrich}{R.~Sennrich, B.~Haddow, A.~Birch, \emph{Neural machine translation of
  rare words with subword units}, arXiv:1508.07909 (2015).
  \url{https://arxiv.org/abs/1508.07909}
  \use{Byte-pair encoding for tokenization.}
  \checked{arXiv record.}}
\refentry{llama3}{A.~Grattafiori, A.~Dubey et al.\ (Llama Team, Meta), \emph{The Llama~3
  herd of models}, arXiv:2407.21783 (2024). \url{https://arxiv.org/abs/2407.21783}
  \use{Concrete numbers for an open model: $405$B parameters, $15.6$T pre-training
  tokens, $128$K-token vocabulary and context, $3.94$ characters per token on English.}
  \checked{arXiv record (current author order starts with Grattafiori); numbers from the
  abstract and the HTML full text.}}
\refentry{kaplan}{J.~Kaplan, S.~McCandlish, T.~Henighan, T.\,B.~Brown et al.,
  \emph{Scaling laws for neural language models}, arXiv:2001.08361 (2020).
  \url{https://arxiv.org/abs/2001.08361}
  \use{Loss as a power law in parameters, data and compute.}
  \checked{arXiv record.}}
\refentry{chinchilla}{J.~Hoffmann, S.~Borgeaud, A.~Mensch, E.~Buchatskaya et al.,
  \emph{Training compute-optimal large language models}, arXiv:2203.15556 (2022).
  \url{https://arxiv.org/abs/2203.15556}
  \use{Model size and data should grow together (``Chinchilla'').}
  \checked{arXiv record.}}
\refentry{christiano}{P.~Christiano, J.~Leike, T.\,B.~Brown, M.~Martic, S.~Legg,
  D.~Amodei, \emph{Deep reinforcement learning from human preferences},
  arXiv:1706.03741 (2017). \url{https://arxiv.org/abs/1706.03741}
  \use{Origin of learning a reward model from human comparisons.}
  \checked{arXiv record.}}
\refentry{instructgpt}{L.~Ouyang, J.~Wu, X.~Jiang, D.~Almeida et al.\ (20 authors),
  \emph{Training language models to follow instructions with human feedback},
  arXiv:2203.02155 (2022). \url{https://arxiv.org/abs/2203.02155}
  \use{RLHF for language models (InstructGPT); 1.3B model preferred over 175B GPT-3.}
  \checked{arXiv record and abstract.}}
\refentry{constitutional}{Y.~Bai, S.~Kadavath, S.~Kundu, A.~Askell et al.,
  \emph{Constitutional AI: Harmlessness from AI feedback}, arXiv:2212.08073 (2022).
  \url{https://arxiv.org/abs/2212.08073}
  \use{Feedback from a model guided by written principles instead of human raters.}
  \checked{arXiv record.}}
\refentry{wei-cot}{J.~Wei, X.~Wang, D.~Schuurmans, M.~Bosma et al.,
  \emph{Chain-of-thought prompting elicits reasoning in large language models},
  arXiv:2201.11903 (2022). \url{https://arxiv.org/abs/2201.11903}
  \use{Intermediate reasoning steps by prompting.}
  \checked{arXiv record.}}
\refentry{tulu3}{N.~Lambert, J.~Morrison, V.~Pyatkin, S.~Huang et al., \emph{Tülu 3:
  Pushing frontiers in open language model post-training}, arXiv:2411.15124 (2024).
  \url{https://arxiv.org/abs/2411.15124}
  \use{Introduces the term ``reinforcement learning with verifiable rewards'' (RLVR).}
  \checked{arXiv record and abstract (``a novel method we call\dots'').}}
\refentry{deepseekr1}{DeepSeek-AI (D.~Guo, D.~Yang, H.~Zhang et al.),
  \emph{DeepSeek-R1: Incentivizing reasoning capability in LLMs via reinforcement
  learning}, arXiv:2501.12948 (2025); published in Nature \textbf{645}, 633 (2025).
  \url{https://arxiv.org/abs/2501.12948}
  \use{Reasoning (self-reflection, verification) emerging from RL on verifiable rewards.}
  \checked{arXiv record, including its journal reference.}}
\refentry{sharma}{M.~Sharma, M.~Tong, T.~Korbak, D.~Duvenaud et al.,
  \emph{Towards understanding sycophancy in language models}, arXiv:2310.13548 (2023).
  \url{https://arxiv.org/abs/2310.13548}
  \use{Sycophancy in five assistants across four tasks; humans and preference models
  sometimes prefer sycophantic answers.}
  \checked{arXiv record and abstract.}}
\refentry{metr-hacking}{S.~Von Arx, L.~Chan, B.~Barnes, \emph{Recent frontier models
  are reward hacking}, METR blog (5 June 2025).
  \url{https://metr.org/blog/2025-06-05-recent-reward-hacking}
  \use{Models modifying tests or scoring code to raise their score.}
  \checked{the page: title, date, authors, quoted phrase.}}
\end{enumerate}

\subsection*{Agents and context}
\begin{enumerate}[label={[\arabic*]}, ref=\arabic*, resume=refs, leftmargin=*, itemsep=6pt]
\refentry{agents}{\best E.~Schluntz, B.~Zhang, \emph{Building effective agents},
  Anthropic engineering blog (19 Dec 2024).
  \url{https://www.anthropic.com/engineering/building-effective-agents}
  \use{Workflows versus agents; start simple.}
  \checked{the page: title and date. \textbf{Partial:} the page names the authors as
  ``Erik S.\ and Barry Zhang''; the surname Schluntz is from secondary sources only.}}
\refentry{react}{S.~Yao, J.~Zhao, D.~Yu, N.~Du et al., \emph{ReAct: Synergizing
  reasoning and acting in language models}, arXiv:2210.03629 (2022).
  \url{https://arxiv.org/abs/2210.03629}
  \use{The reason--act--observe loop behind agents.}
  \checked{arXiv record.}}
\refentry{swebench}{C.\,E.~Jimenez, J.~Yang, A.~Wettig, S.~Yao et al., \emph{SWE-bench:
  Can language models resolve real-world GitHub issues?}, arXiv:2310.06770 (2023).
  \url{https://arxiv.org/abs/2310.06770}
  \use{2294 real issues from 12 Python repositories, judged by their own tests.}
  \checked{arXiv record and abstract.}}
\refentry{metr-horizon}{T.~Kwa, B.~West, J.~Becker, A.~Deng et al., \emph{Measuring AI
  ability to complete long software tasks}, arXiv:2503.14499 (2025), NeurIPS 2025.
  \url{https://arxiv.org/abs/2503.14499}
  \use{The $50\,\%$ time horizon; doubling about every seven months since 2019.}
  \checked{arXiv record and abstract.}}
\refentry{contextrot}{\best K.~Hong, A.~Troynikov, J.~Huber, \emph{Context rot: How
  increasing input tokens impacts LLM performance}, Chroma technical report (14 July
  2025). \url{https://www.trychroma.com/research/context-rot}
  \use{18 models; degradation with input length; distractors; focused (${\sim}300$
  tokens) versus full (${\sim}113$K) prompts; shuffled haystacks easier.}
  \checked{the report page (old URL \cmd{research.trychroma.com} redirects here).}}
\refentry{lostmiddle}{N.\,F.~Liu, K.~Lin, J.~Hewitt, A.~Paranjape et al., \emph{Lost in
  the middle: How language models use long contexts}, arXiv:2307.03172 (2023).
  \url{https://arxiv.org/abs/2307.03172}
  \use{Information in the middle of a long context is used worst.}
  \checked{arXiv record.}}
\refentry{context-eng}{\best P.~Rajasekaran, E.~Dixon, C.~Ryan, J.~Hadfield (Anthropic
  Applied AI team), \emph{Effective context engineering for AI agents}, Anthropic
  engineering blog (29 Sep 2025).
  \url{https://www.anthropic.com/engineering/effective-context-engineering-for-ai-agents}
  \use{Context as a finite resource; compaction; note-taking; sub-agents; just-in-time
  retrieval.}
  \checked{the page: title, date, authors.}}
\refentry{cc-best}{Anthropic, \emph{Best practices for Claude Code}, Claude Code
  documentation (living document, no date or author given).
  \url{https://code.claude.com/docs/en/best-practices}
  \use{Tool-specific, but its core is general: give the agent a way to verify its work;
  explore, plan, implement, commit; keep instruction files short; clear context between
  tasks.}
  \checked{the page. The former URL
  \cmd{anthropic.com/engineering/claude-code-best-practices} (2025 blog post) now
  redirects here.}}
\refentry{agentsmd}{\emph{AGENTS.md}, an open format for guiding coding agents, stewarded
  by the Agentic AI Foundation (Linux Foundation). \url{https://agents.md/}
  \use{Which agents read \cmd{AGENTS.md}.}
  \checked{the page.}}
\refentry{skills}{B.~Zhang, K.~Lazuka, M.~Murag, \emph{Equipping agents for the real
  world with Agent Skills}, Anthropic engineering blog (16 Oct 2025).
  \url{https://www.anthropic.com/engineering/equipping-agents-for-the-real-world-with-agent-skills}
  \use{Skills and progressive disclosure.}
  \checked{the page: title, date, authors.}}
\refentry{mcp}{Anthropic, \emph{Introducing the Model Context Protocol} (25 Nov 2024).
  \url{https://www.anthropic.com/news/model-context-protocol}; specification at
  \url{https://modelcontextprotocol.io/}
  \use{Trends section.}
  \checked{both pages.}}
\refentry{willison}{S.~Willison, \emph{Simon Willison's Weblog}.
  \url{https://simonwillison.net/}
  \use{Frequent, careful, hands-on notes on LLM tools; good for staying current.}
  \checked{the site.}}
\end{enumerate}

\subsection*{Research software engineering}
\begin{enumerate}[label={[\arabic*]}, ref=\arabic*, resume=refs, leftmargin=*, itemsep=6pt]
\refentry{wilson-good}{\best G.~Wilson, J.~Bryan, K.~Cranston, J.~Kitzes,
  L.~Nederbragt, T.\,K.~Teal, \emph{Good enough practices in scientific computing},
  PLOS Comput.\ Biol.\ \textbf{13}, e1005510 (2017).
  \url{https://doi.org/10.1371/journal.pcbi.1005510}
  \use{The minimum set of practices for research code; short and practical.}
  \checked{Crossref.}}
\refentry{wilson-best}{G.~Wilson, D.\,A.~Aruliah, C.\,T.~Brown, N.\,P.~Chue Hong et al.,
  \emph{Best practices for scientific computing}, PLOS Biol.\ \textbf{12}, e1001745
  (2014). \url{https://doi.org/10.1371/journal.pbio.1001745}
  \use{The more ambitious predecessor.}
  \checked{Crossref.}}
\refentry{turingway}{The Turing Way Community, \emph{The Turing Way: a handbook for
  reproducible, ethical and collaborative data science}.
  \url{https://book.the-turing-way.org/}
  \use{Reference handbook on reproducibility, version control, testing.}
  \checked{the site; the old address \cmd{the-turing-way.netlify.app} still works.}}
\refentry{progit}{\best S.~Chacon, B.~Straub, \emph{Pro Git}, 2nd ed.\ (Apress, 2014),
  free online. \url{https://git-scm.com/book/en/v2}
  \use{Chapters 2 (basics), 3 (branching), 10 (internals).}
  \checked{the site.}}
\refentry{swc-git}{\best Software Carpentry, \emph{Version Control with Git}.
  \url{https://swcarpentry.github.io/git-novice/}
  \use{Self-paced git lesson for researchers.}
  \checked{the site.}}
\refentry{swc-testing}{The Carpentries Incubator, \emph{Python Testing and Continuous
  Integration} (alpha). \url{https://carpentries-incubator.github.io/python-testing/}
  \use{Testing and CI in Python; the ideas transfer to Julia.}
  \checked{the site. Note: an incubator lesson, not an official Software Carpentry one.}}
\refentry{scipy-dev}{Scientific Python, \emph{Scientific Python Library Development
  Guide}. \url{https://learn.scientific-python.org/development/}
  \use{Packaging, testing and CI for Python research code.}
  \checked{the site.}}
\refentry{drwatson}{G.~Datseris, J.~Isensee, S.~Pech, T.~Gál, \emph{DrWatson: the
  perfect sidekick for your scientific inquiries}, J.\ Open Source Softw.\ \textbf{5},
  2673 (2020). \url{https://doi.org/10.21105/joss.02673}
  \use{Citation for DrWatson.}
  \checked{Crossref and the JOSS page.}}
\refentry{drwatson-docs}{\best DrWatson.jl documentation.
  \url{https://juliadynamics.github.io/DrWatson.jl/stable/}
  \use{\cmd{produce\_or\_load}, \cmd{tagsave}, \cmd{savename}; the ``Workflow
  tutorial''.}
  \checked{the pages on saving and naming; \cmd{savename} and \cmd{dict\_list} output
  also checked by running DrWatson.}}
\refentry{julia}{J.~Bezanson, A.~Edelman, S.~Karpinski, V.\,B.~Shah, \emph{Julia: A fresh
  approach to numerical computing}, SIAM Rev.\ \textbf{59}, 65 (2017).
  \url{https://doi.org/10.1137/141000671}
  \use{Citation for Julia.}
  \checked{Crossref.}}
\refentry{julia-test}{Julia documentation, \emph{Unit Testing} (standard library
  \cmd{Test}). \url{https://docs.julialang.org/en/v1/stdlib/Test/}
  \use{\cmd{@test}, \cmd{@testset}; \cmd{isapprox} tolerances (behaviour checked by
  running Julia 1.13).}
  \checked{the page.}}
\refentry{julia-actions}{julia-actions, GitHub Actions for Julia.
  \url{https://github.com/julia-actions}
  \use{\cmd{setup-julia}, \cmd{cache}, \cmd{julia-buildpkg}, \cmd{julia-runtest}.}
  \checked{the organisation page and the \cmd{julia-runtest} repository.}}
\refentry{matlab-unittest}{MathWorks, \emph{matlab.unittest.TestCase} (MATLAB
  documentation).
  \url{https://www.mathworks.com/help/matlab/ref/matlab.unittest.testcase-class.html}
  \use{Unit tests in MATLAB.}
  \checked{the page.}}
\end{enumerate}

\subsection*{Physics of the toy problem}
\begin{enumerate}[label={[\arabic*]}, ref=\arabic*, resume=refs, leftmargin=*, itemsep=6pt]
\refentry{krantz}{\best P.~Krantz, M.~Kjaergaard, F.~Yan, T.\,P.~Orlando,
  S.~Gustavsson, W.\,D.~Oliver, \emph{A quantum engineer's guide to superconducting
  qubits}, Appl.\ Phys.\ Rev.\ \textbf{6}, 021318 (2019); arXiv:1904.06560.
  \url{https://doi.org/10.1063/1.5089550}
  \use{Rotating frame, Rabi, Ramsey, echo, $T_1$, $T_2$, from an experimental viewpoint.}
  \checked{Crossref and arXiv record.}}
\refentry{rabi}{I.\,I.~Rabi, \emph{Space quantization in a gyrating magnetic field},
  Phys.\ Rev.\ \textbf{51}, 652 (1937). \url{https://doi.org/10.1103/PhysRev.51.652}
  \use{Historical origin of the Rabi formula.}
  \checked{Crossref.}}
\refentry{lindblad}{G.~Lindblad, \emph{On the generators of quantum dynamical
  semigroups}, Commun.\ Math.\ Phys.\ \textbf{48}, 119 (1976).
  \url{https://doi.org/10.1007/BF01608499}
  \checked{Crossref.}}
\refentry{gks}{V.~Gorini, A.~Kossakowski, E.\,C.\,G.~Sudarshan, \emph{Completely positive
  dynamical semigroups of $N$-level systems}, J.\ Math.\ Phys.\ \textbf{17}, 821 (1976).
  \url{https://doi.org/10.1063/1.522979}
  \use{With [\ref{ref:lindblad}]: the general form of a Markovian master equation.}
  \checked{Crossref.}}
\refentry{breuer}{H.-P.~Breuer, F.~Petruccione, \emph{The Theory of Open Quantum
  Systems} (Oxford University Press, 2002).
  \url{https://doi.org/10.1093/acprof:oso/9780199213900.001.0001}
  \use{Textbook derivation of the master equation.}
  \checked{Crossref (the DOI is the 2007 online edition).}}
\refentry{ramsey}{N.\,F.~Ramsey, \emph{A molecular beam resonance method with separated
  oscillating fields}, Phys.\ Rev.\ \textbf{78}, 695 (1950).
  \url{https://doi.org/10.1103/PhysRev.78.695}
  \checked{Crossref.}}
\refentry{hahn}{E.\,L.~Hahn, \emph{Spin echoes}, Phys.\ Rev.\ \textbf{80}, 580 (1950).
  \url{https://doi.org/10.1103/PhysRev.80.580}
  \checked{Crossref.}}
\refentry{jaynescummings}{E.\,T.~Jaynes, F.\,W.~Cummings, \emph{Comparison of quantum and
  semiclassical radiation theories with application to the beam maser},
  Proc.\ IEEE \textbf{51}, 89 (1963). \url{https://doi.org/10.1109/PROC.1963.1664}
  \checked{Crossref.}}
\refentry{qojulia}{S.~Krämer, D.~Plankensteiner, L.~Ostermann, H.~Ritsch,
  \emph{QuantumOptics.jl: A Julia framework for simulating open quantum systems},
  Comput.\ Phys.\ Commun.\ \textbf{227}, 109 (2018).
  \url{https://doi.org/10.1016/j.cpc.2018.02.004}
  \use{The independent route of block~6.}
  \checked{Crossref.}}
\refentry{qojulia-docs}{QuantumOptics.jl documentation. \url{https://docs.qojulia.org/}
  \use{\cmd{timeevolution.master(tspan, rho0, H, J)}; \cmd{SpinBasis(1//2)} with
  \cmd{spinup} $= \ket{e}$.}
  \checked{the page on master equations; basis convention and solver checked by running
  QuantumOptics.jl 1.2.}}
\refentry{lsqfit}{LsqFit.jl documentation.
  \url{https://julianlsolvers.github.io/LsqFit.jl/latest/}
  \use{\cmd{curve\_fit} with inverse-variance weights; \cmd{coef}, \cmd{stderror}.}
  \checked{the tutorial page; function names checked by running LsqFit 0.16.}}
\end{enumerate}

\subsection*{Responsibility and policy}
\begin{enumerate}[label={[\arabic*]}, ref=\arabic*, resume=refs, leftmargin=*, itemsep=6pt]
\refentry{nature-rules}{Editorial, \emph{Tools such as ChatGPT threaten transparent
  science; here are our ground rules for their use}, Nature \textbf{613}, 612 (2023).
  \url{https://doi.org/10.1038/d41586-023-00191-1}
  \use{No AI author; document use in methods or acknowledgements.}
  \checked{Crossref. \textbf{Partial:} the publisher page required a login; the two
  rules are from consistent secondary reproductions.}}
\refentry{aps-news}{American Physical Society, \emph{American Physical Society releases
  updated AI policy for APS journals} (17 June 2026).
  \url{https://www.aps.org/about/news/2026/06/releases-updated-ai-policy-journals}
  \checked{the page.}}
\refentry{aps-policy}{\best APS Journals, \emph{Appropriate Use of AI Tools}.
  \url{https://journals.aps.org/authors/appropriate-use-ai-tools}
  \use{Current rules for authors and referees of APS journals.}
  \checked{the page (undated).}}
\end{enumerate}

\subsection*{Trends}
\begin{enumerate}[label={[\arabic*]}, ref=\arabic*, resume=refs, leftmargin=*, itemsep=6pt]
\refentry{jevons}{W.\,S.~Jevons, \emph{The Coal Question} (Macmillan, London, 1865;
  2nd ed.\ 1866). Full text of the 2nd edition:
  \url{https://www.econlib.org/library/YPDBooks/Jevons/jvnCQ.html}
  \use{The Jevons paradox; Chapter~VII, ``Of the Economy of Fuel''.}
  \checked{the Econlib text (chapter title and quoted sentence).}}
\refentry{vibecoding}{A.~Karpathy, post on X introducing ``vibe coding'' (2 Feb 2025).
  \url{https://x.com/karpathy/status/1886192184808149383}
  \checked{\textbf{Partial:} the post did not render without login; date and wording
  from several consistent secondary sources.}}
\end{enumerate}

\subsection*{Claims on the slides without a checkable source}
\begin{itemize}
  \item ``About $0.75$ words per token'' (slide, block~1) is a widely quoted rule of thumb
        from OpenAI's help pages; the page refused automated access. Section~\ref{sec:llm}
        supports the number with the Llama~3 figure of $3.94$ characters per token
        \rf{llama3} instead.
\end{itemize}


\end{document}
